
\setcounter{chapter}{17}

\chapter{表示论与特征标理论}

\section{线性作用与群环上的模}

本书余下部分中，除非另有说明，我们考虑的群都是有限群。本节中 \(F\) 始终是域，\(G\) 是有限群。我们先引入基本术语。回忆若 \(V\) 是 \(F\) 上向量空间，则 \(GL(V)\) 是 \(V\) 到自身的非奇异线性变换群（在复合下），而若 \(n \in \mathbb{Z}^+\)，则 \(GL_n(F)\) 是元素取自 \(F\) 的可逆 \(n \times n\) 矩阵群（在矩阵乘法下）。

\begin{definition}
设 \(G\) 是有限群，\(F\) 是域，\(V\) 是 \(F\) 上向量空间。

（1）\(G\) 的\textbf{线性表示}是任何从 \(G\) 到 \(GL(V)\) 的同态。表示的\textbf{次数}是 \(V\) 的维数。

（2）设 \(n \in \mathbb{Z}^+\). \(G\) 的\textbf{矩阵表示}是任何从 \(G\) 到 \(GL_n(F)\) 的同态。

（3）若线性或矩阵表示是单射，则称它是\textbf{忠实的}。

（4）\(G\) 在 \(F\) 上的\textbf{群环}是形如 \(\sum_{g \in G}a_gg\) 的所有形式和之集，其加法按分量进行，乘法为 \((ag)(\beta h) = (a\beta)(gh)\)（其中 \(a\) 与 \(\beta\) 在 \(F\) 中相乘，\(gh\) 是 \(G\) 中的乘积），并按分配律延拓到和上（参见第 7.2 节）。
\end{definition}

除非我们专门讨论置换表示，术语「表示」总是指「线性表示」。当我们想强调域 \(F\) 时，我们说 \(F\)-表示，或者说 \(G\) 在 \(V\) 上关于 \(F\) 的表示。

回忆若 \(V\) 是 \(n\) 维有限维向量空间，则通过固定 \(V\) 的一组基，我们得到同构 \(GL(V) \cong GL_n(F)\). 这样 \(G\) 在有限维向量空间上的任何线性表示都给出一个矩阵表示，反之亦然。我们的线性表示大多是有限次的，我们将在线性表示与矩阵表示之间自由切换（当想给出两者的显式对应时指定一组基）。此外，给定有限次的线性表示 \(\varphi : G \to GL(V)\)，相应的矩阵表示提供了一些数值不变量（例如对 \(g \in G\) 的 \(\det\varphi(g)\)），它们不依赖于给出 \(GL(V)\) 与 \(GL_n(F)\) 之间同构的基的选取。利用这样的不变量将是我们展开的基础。

在给出表示的例子之前，我们更详细地回忆群环 \(FG\)（群环在第 7.2 节引入，该节讨论了一些记号与例子）。设 \(G\) 的元素为 \(g_1, g_2, \ldots, g_n\). \(FG\) 的每个元素都具有形式
\[
\sum_{i=1}^{n}a_ig_i, \qquad a_i \in F.
\]
两个形式和相等当且仅当所有相应群元素的系数相等。\(FG\) 中的加法与乘法定义如下：
\[
\sum_{i=1}^{n}a_ig_i+\sum_{i=1}^{n}\beta_ig_i = \sum_{i=1}^{n}(a_i+\beta_i)g_i, \qquad \left(\sum_{i=1}^{n}a_ig_i\right)\left(\sum_{j=1}^{n}\beta_jg_j\right) = \sum_{k}\left(\sum_{g_ig_j = g_k}a_i\beta_j\right)g_k,
\]
其中系数 \(a_i\) 与 \(\beta_i\) 的加法与乘法在 \(F\) 中进行。注意按乘法的定义，

\noindent\textbf{\(FG\) 是交换环当且仅当 \(G\) 是阿贝尔群。}

群 \(G\) 出现在 \(FG\) 中（把 \(g_i\) 与 \(1g_i\) 等同），域 \(F\) 出现在 \(FG\) 中（把 \(\beta\) 与 \(\beta g_1\) 等同，其中 \(g_1\) 是 \(G\) 的单位元）。在这些等同下
\[
\beta\left(\sum_{i=1}^{n}a_ig_i\right) = \sum_{i=1}^{n}(\beta a_i)g_i, \qquad \text{对所有 } \beta \in F.
\]
这样

\noindent\textbf{\(FG\) 是 \(F\) 上以 \(G\) 的元素为基的向量空间。}

特别地，\(FG\) 是 \(F\) 上维数为 \(|G|\) 的向量空间。\(F\) 的元素与 \(FG\) 的所有元素交换，即 \(F\) 在 \(FG\) 的中心中。当我们想强调后两条性质时，我们说 \(FG\) 是 \(F\)-\textbf{代数}（一般地，\(F\)-代数是其中心含有 \(F\) 的环 \(R\)，故 \(R\) 既是环又是 \(F\)-向量空间）。

注意 \(FG\) 中的运算类似于 \(F\)-代数 \(F[x]\) 中的运算（尽管 \(F[x]\) 在 \(F\) 上是无限维的）。在某些著作中 \(FG\) 记作 \(F[G]\)，尽管后者目前不常见。

\noindent\textbf{例}

（1）若 \(G = \langle g\rangle\) 是 \(n \in \mathbb{Z}^+\) 阶循环群，则 \(FG\) 的元素具有形式 \(\sum_{i=0}^{n-1}a_ig^i\). 把 \(x^k\) 映到 \(g^k\)（\(k \geq 0\)）的映射 \(F[x] \to F\langle g\rangle\) 按 \(F\)-线性延拓为核等于由 \(x^n-1\) 生成的理想 \((x^n-1)\) 的满环同态。故
\[
F\langle g\rangle \cong F[x]/(x^n-1).
\]
这是 \(F\)-代数的同构，即是 \(F\)-线性的环同构。

（2）用上面例 1 的记号，设 \(g_1 + g + g^2+\cdots+g^{n-1}\) 记为 \(N\)，\(r = 1-g\). 则 \(N(1-g) = 0\). 故环 \(F\langle g\rangle\) 含有零因子（只要 \(n > 1\)）。更一般地，若 \(G\) 是阶大于 1 的任意群，则对任意非单位元 \(g \in G\)，\(F\langle g\rangle\) 是 \(FG\) 的子环，故 \(FG\) 也含有零因子。

（3）设 \(G = S_3\)、\(F = \mathbb{Q}\). 元素 \(r = 5(1\ 2)-7(1\ 2\ 3)\) 与 \(s = -4(1\ 2\ 3)+12(1\ 3\ 2)\) 是 \(\mathbb{Q}S_3\) 的典型成员。它们的和与积为
\[
r+s = 5(1\ 2)-11(1\ 2\ 3)+12(1\ 3\ 2),
\]
\[
rs = -20(2\ 3)+28(1\ 3\ 2)+60(1\ 3)-84
\]
（回忆置换的乘积（复合）从右往左计算）。群环 \(\mathbb{Q}D_8\) 中两个元素的和与积的显式例子见第 7.2 节。

在给出表示的具体例子之前，我们讨论 \(G\) 的表示与 \(FG\)-模之间的对应（之后可以同时给出两者的例子）。这段讨论与第 10.1 节中对 \(F[x]\)-模的处理非常平行。

先设 \(\varphi : G \to GL(V)\) 是 \(G\) 在 \(F\) 上向量空间 \(V\) 上的表示。取一个 \(F\)-基 \(\{v_1, \ldots, v_n\}\)，并对每个 \(i \in \{1, \ldots, n\}\)，通过令 \(g_i \cdot v_j\) 为 \(\varphi(g_i)\) 在 \(v_j\) 上的作用，把 \(V\) 变成一个 \(FG\)-模，即定义一个环元素对 \(V\) 的元素的作用如下：
\[
\left(\sum_{i=1}^{n}a_ig_i\right)\cdot v = \sum_{i=1}^{n}a_i(g_i \cdot v) \qquad \text{对所有 } \sum_{i=1}^{n}a_ig_i \in FG,\ v \in V.
\]
我们验证模公理 2（b）的一个特殊情形，它恰好表明「\(\varphi\) 是群同态」这一事实在哪里被用到：
\[
(g_ig_j)\cdot v = \varphi(g_ig_j)(v) = (\varphi(g_i)\circ\varphi(g_j))(v) = \varphi(g_i)(\varphi(g_j)(v)) = g_i \cdot (g_j \cdot v).
\]
这个论证按线性性可延拓到 \(FG\) 的任意元素，从而证明模公理 2（b）一般成立。其余模公理的验证是一个习题。

注意 \(F\) 是 \(FG\) 的子环，域元素 \(a\) 对向量的作用与环元素 \(a\) 对向量的作用相同，即 \(FG\)-模作用延拓了 \(V\) 上的 \(F\)-作用。

现在反之设给定一个 \(FG\)-模 \(V\). 我们如下得到相应的 \(F\) 上向量空间与 \(G\) 的表示。由于 \(V\) 是 \(FG\)-模，它是 \(F\)-模，即 \(F\) 上向量空间。又对每个 \(g \in G\) 我们得到从 \(V\) 到 \(V\) 的映射，记作 \(\varphi(g)\)，定义为
\[
\varphi(g)(v) = g \cdot v \qquad \text{对所有 } v \in V,
\]
其中 \(g \cdot v\) 是给定的环元素 \(g\) 对 \(V\) 的元素 \(v\) 的作用。由于 \(F\) 的元素与每个 \(g \in G\) 交换，由模公理推出对所有 \(v, w \in V\) 与所有 \(a, \beta \in F\) 有
\[
\varphi(g)(av+\beta w) = g \cdot (av+\beta w) = g \cdot (av)+g \cdot (\beta w) = a(g \cdot v)+\beta(g \cdot w) = a\varphi(g)(v)+\beta\varphi(g)(w),
\]
即对每个 \(g \in G\)，\(\varphi(g)\) 是线性变换。进一步，由模公理 2（b）推出 \(\varphi(gh) = \varphi(g)\varphi(h)\)（这本质上就是把上面的计算反过来）。这证明 \(\varphi\) 是群同态（特别地 \(\varphi(g^{-1}) = \varphi(g)^{-1}\)，故 \(G\) 的每个元素映到非奇异线性变换，即 \(\varphi : G \to GL(V)\)）。

这段讨论表明 \(FG\)-模与「\(V\) 是 \(F\) 上向量空间连同表示 \(\varphi : G \to GL(V)\)」这样的对之间存在双射。因此给出 \(V\) 上的表示 \(\varphi : G \to GL(V)\) 等价于给出一个 \(FG\)-模 \(V\). 在这个对应下我们说模 \(V\) \textbf{提供}（affords）\(G\) 的表示 \(\varphi\).

回忆第 10.1 节：若通过线性变换 \(T\) 把向量空间 \(M\) 变成 \(F[x]\)-模，则 \(M\) 的 \(F[x]\)-子模恰是 \(M\) 的 \(T\)-不变子空间。在当前情形，若 \(V\) 是提供表示 \(\varphi\) 的 \(FG\)-模，则 \(V\) 的子空间 \(U\) 称为 \textbf{\(G\)-不变}或 \textbf{\(G\)-稳定}的，若对所有 \(g \in G\) 与所有 \(u \in U\) 有 \(g \cdot u \in U\)（即若对所有 \(g \in G\) 与所有 \(u \in U\) 有 \(\varphi(g)(u) \in U\)）。由此容易推出

\noindent\textbf{\(V\) 的 \(FG\)-子模恰是 \(V\) 的 \(G\)-稳定子空间。}

\noindent\textbf{例}

（1）设 \(V\) 是 \(F\) 上 1 维向量空间，通过令对所有 \(g \in G\) 与 \(v \in V\) 有 \(gv = v\) 把 \(V\) 变成 \(FG\)-模。这个模提供由 \(\varphi(g) = I\)（恒等线性变换，对所有 \(g \in G\)）定义的表示 \(\varphi : G \to GL(V)\). 相应的矩阵表示（关于 \(V\) 的任意基）是把每个群元素映到 \(1 \times 1\) 单位矩阵的同态 \(G \to GL_1(F)\). 今后我们把它称为 \(G\) 的\textbf{平凡表示}。平凡表示的次数为 1，且若 \(|G| > 1\) 则它不是忠实的。

（2）设 \(V = FG\) 并把这个环看作自身的左模。则 \(V\) 提供 \(G\) 的次数等于 \(|G|\) 的表示。若取 \(G\) 的元素作为 \(V\) 的基，则每个 \(g \in G\) 在左正则置换表示下置换这些基元素：\(g \cdot g_i = gg_i\). 关于 \(V\) 的这组基，群元素 \(g\) 的矩阵在 \(gg_j = g_i\) 时第 \(i\) 行第 \(j\) 列为 1，其他位置为 0. 这个（线性或矩阵）表示称为 \(G\) 的\textbf{正则表示}。注意 \(G\) 的每个非单位元诱导 \(V\) 的基上的非单位置换，故正则表示总是忠实的。

（3）设 \(n \in \mathbb{Z}^+\)，\(G = S_n\)，\(V\) 是 \(F\) 上以 \(e_1, e_2, \ldots, e_n\) 为基的 \(n\) 维向量空间。让 \(S_n\) 作用在 \(V\) 上：对每个 \(\sigma \in S_n\) 定义 \(\sigma \cdot e_i = e_{\sigma(i)}\)，即 \(\sigma\) 通过置换基元素的下标作用。这给出 \(S_n\) 到 \(GL(V)\) 的（单）同态（即 \(S_n\) 的忠实的 \(n\) 次表示），从而把 \(V\) 变成 \(FS_n\)-模。如上一例，关于基 \(e_1, \ldots, e_n\)，\(\sigma\) 的矩阵在 \(\sigma \cdot e_j = e_i\) 时第 \(i\) 行第 \(j\) 列为 1（其他位置为 0）。故当 \(\sigma(j) = i\) 时 \(\sigma\) 的矩阵第 \(i\) 行第 \(j\) 列为 1.

作为环作用的一个例子，考虑 \(FS_3\) 在以 \(e_1, e_2, e_3\) 为基的 \(F\) 上 3 维向量空间上的作用。设 \(\sigma\) 是对换 \((1\ 2)\)，\(\tau\) 是 3-轮换 \((1\ 2\ 3)\)，\(r = 2\sigma-3\tau \in FS_3\). 则
\[
r \cdot (ae_1+\beta e_2+\gamma e_3) = 2(ae_{\sigma(1)}+\beta e_{\sigma(2)}+\gamma e_{\sigma(3)})-3(ae_{\tau(1)}+\beta e_{\tau(2)}+\gamma e_{\tau(3)})
\]
\[
= 2(ae_2+\beta e_1+\gamma e_3)-3(ae_2+\beta e_3+\gamma e_1) = (2\beta-3\gamma)e_1-ae_2+(2\gamma-3\beta)e_3.
\]

（4）若 \(\psi : H \to GL(V)\) 是 \(H\) 的任意表示，\(\varphi : G \to H\) 是任意群同态，则复合 \(\psi \circ \varphi\) 是 \(G\) 的表示。例如设 \(V\) 是上一例中描述的 \(n\) 维 \(FS_n\)-模。若 \(\pi : G \to S_n\) 是 \(G\) 的任意置换表示，\(\pi\) 与上述表示的复合给出 \(G\) 的线性表示。换句话说，\(V\) 在作用 \(g \cdot e_i = e_{\pi(g)(i)}\)（对所有 \(g \in G\)）下成为 \(FG\)-模。注意正则表示（2）就是这里 \(\pi\) 为 \(G\) 的左正则置换表示、\(n = |G|\) 的特殊情形。

（5）任何从 \(G\) 到乘法群 \(F^\times = GL_1(F)\) 的同态都是次数为 1 的（矩阵）表示。例如设 \(G = \langle g\rangle \cong \mathbb{Z}_n\) 是 \(n\) 阶循环群，\(s\) 是 \(F\) 中固定的 \(n\) 次单位根。令 \(g^i \mapsto s^i\)（对所有 \(i \in \mathbb{Z}\)）。这个 \(\langle g\rangle\) 的表示是忠实的当且仅当 \(s\) 是本原 \(n\) 次单位根。

（6）在许多情形给出群 \(G\) 的显式矩阵表示比给出 \(FG\)-模更容易。例如回忆二面体群 \(D_{2n}\) 有表现
\[
D_{2n} = \langle r, s \mid r^n = s^2 = 1,\ rs = sr^{-1}\rangle.
\]
若 \(R\) 与 \(S\) 是满足关系 \(R^n = S^2 = I\) 与 \(RS = SR^{-1}\) 的任意矩阵，则映射 \(r \mapsto R\)、\(s \mapsto S\) 唯一地延拓为从 \(D_{2n}\) 到由 \(R\) 与 \(S\) 生成的矩阵群的同态，从而给出 \(D_{2n}\) 的一个表示。\(M_2(\mathbb{R})\) 中矩阵 \(R, S\) 的一个显式例子可以如下得到。若在 \(x,y\) 平面上以原点为中心画一个正 \(n\) 边形并使其一条对称轴是直线 \(y = x\)，则把平面旋转 \(2\pi/n\) 弧度的矩阵 \(R\) 与把平面关于直线 \(y = x\) 反射的矩阵 \(S\) 都把这个 \(n\) 边形映到自身。由此这些矩阵作用为该 \(n\) 边形的对称，从而满足上述关系。这些矩阵容易算出（参见第 1.6 节习题 25），故映射
\[
r \mapsto R = \begin{pmatrix}\cos\frac{2\pi}{n} & -\sin\frac{2\pi}{n}\\ \sin\frac{2\pi}{n} & \cos\frac{2\pi}{n}\end{pmatrix}, \qquad s \mapsto S = \begin{pmatrix}0 & 1\\ 1 & 0\end{pmatrix}
\]
唯一地延拓为 \(D_{2n}\) 到 \(GL_2(\mathbb{R})\) 的（2 次）表示。由于矩阵 \(R\) 与 \(S\) 的阶分别为 \(n\) 与 2，它们生成 \(GL_2(\mathbb{R})\) 的 \(2n\) 阶子群，故此表示是忠实的。

（7）利用四元数群的通常生成元与关系，可以类似地（参见第 1.6 节习题 26）得到从 \(Q_8\) 到 \(GL_2(\mathbb{C})\) 的表示 \(\varphi\)，由
\[
\varphi(i) = \begin{pmatrix}\sqrt{-1} & 0\\ 0 & -\sqrt{-1}\end{pmatrix}, \qquad \varphi(j) = \begin{pmatrix}0 & -1\\ 1 & 0\end{pmatrix}
\]
定义。这个 \(Q_8\) 的表示是忠实的。

（8）四元数群 \(Q_8\) 的 4 维表示可以由实 Hamilton 四元数得到（参见第 7.1 节）。群 \(Q_8\) 是 \(\mathbb{H}\) 的单位四元数乘法群的子群，而 \(Q_8\) 的每个元素通过左乘作用在 4 维实向量空间 \(\mathbb{H}\) 上。由于实数在 \(\mathbb{H}\) 的中心中（即 \(\mathbb{H}\) 是 \(\mathbb{R}\)-代数），左乘是 \(\mathbb{R}\)-线性的。这个线性作用于是给出从 \(Q_8\) 到 \(GL_4(\mathbb{R})\) 的同态。容易写出 \(Q_8\) 的每个元素关于基 \(1, i, j, k\) 的显式矩阵。例如左乘 \(i\) 的作用为 \(1 \mapsto i\)、\(i \mapsto -1\)、\(j \mapsto k\)、\(k \mapsto -j\)，而左乘 \(j\) 的作用为 \(1 \mapsto j\)、\(i \mapsto -k\)、\(j \mapsto -1\)、\(k \mapsto i\). 这个 \(Q_8\) 的表示也是忠实的。

（9）设 \(H\) 是群 \(G\) 的正规子群，且 \(H\) 是某个素数 \(p\) 的初等阿贝尔 \(p\)-群。则 \(V = H\) 是 \(\mathbb{F}_p\) 上的向量空间，其中纯量 \(a\) 作用在向量 \(v\) 上为 \(av = v^a\)（参见第 10.1 节）。\(G\) 的每个元素通过共轭作用在 \(V\) 上是 \(\mathbb{F}_p\)-线性的，因为 \(gv^ag^{-1} = (gvg^{-1})^a\)，而 \(G\) 在 \(V\) 上的这个作用把 \(V\) 变成 \(\mathbb{F}_pG\)-模（初等阿贝尔 \(p\)-群的自同构在第 4.4 与 10.1 节中讨论过）。这个表示的核是与 \(H\) 的每个元素交换的 \(G\) 的元素之集 \(C_G(H)\)（它总含有阿贝尔群 \(H\) 本身）。故群对自身子集的作用常常提供有限域上的线性表示。群在有限域上的表示称为\textbf{模表示}，它们对群内部结构的研究是基本的。

（10）作为 \(FG\)-子模的例子，设 \(G = S_n\)，\(V\) 是例 3 中描述的 \(FS_n\)-模。设 \(N\) 是由所有坐标都相等的向量构成的 \(V\) 的子空间，即
\[
N = \{a_1e_1+a_2e_2+\cdots+a_ne_n \mid a_1 = a_2 = \cdots = a_n\}
\]
（这是 1 维 \(S_n\)-稳定子空间）。每个 \(\sigma \in S_n\) 固定 \(N\) 中每个向量，故子模 \(N\) 提供 \(S_n\) 的平凡表示。作为一个习题，可以证明若 \(n \geq 3\) 则 \(N\) 是 \(V\) 中唯一的 1 维 \(S_n\)-稳定子空间，即 \(N\) 是唯一的 1 维 \(FS_n\)-子模（\(N\) 称为 \(FS_n\) 的\textbf{迹子模}）。

\(V\) 的另一个 \(FS_n\)-子模是坐标和为零的所有向量构成的子空间 \(I\)：
\[
I = \{a_1e_1+a_2e_2+\cdots+a_ne_n \mid a_1+a_2+\cdots+a_n = 0\}.
\]
同样 \(I\) 是 \(S_n\)-稳定子空间（因为每个 \(\sigma \in S_n\) 置换 \(V\) 中每个向量的坐标，从而保持系数之和不变）。由于 \(I\) 是把向量映到其系数之和的线性变换 \(V \to F\)（称为增广映射，参见第 7.3 节）的核，\(I\) 的维数为 \(n-1\).

（11）若 \(V = FG\) 是上面例 2 中描述的 \(G\) 的正则表示，则如上一例 \(V\) 有维数为 1 与 \(|G|-1\) 的 \(FG\)-子模：
\[
N = \left\{\sum_{i=1}^{n}a_ig_i \;\middle|\; a_1 = a_2 = \cdots = a_n\right\}, \qquad I = \left\{\sum_{i=1}^{n}a_ig_i \;\middle|\; a_1+a_2+\cdots+a_n = 0\right\}.
\]
事实上 \(N\) 与 \(I\) 是 \(FG\) 的双边理想（不只是左理想——注意 \(N\) 在 \(FG\) 的中心中）。理想 \(I\) 称为 \(FG\) 的\textbf{增广理想}，\(N\) 称为 \(FG\) 的\textbf{迹理想}。

回忆在研究向量空间 \(V\) 到自身的线性变换 \(T\) 时，我们把 \(V\) 变成 \(F[x]\)-模（其中 \(x\) 在 \(V\) 上作用为 \(T\)）；我们的目的是把 \(V\) 分解为循环子模的直和。这样我们能够找到 \(V\) 的一组基，使 \(T\) 关于这组基的矩阵具有某种标准形。改变 \(V\) 的基不改变模 \(V\)，但通过相似改变 \(T\) 的矩阵表示（即改变 \(GL(V)\) 与 \(GL_n(F)\) 之间的同构）。我们引入类似的术语来描述两个 \(FG\)-模在相差基变换的意义下相同。

\begin{definition}
若提供两个表示的 \(FG\)-模是同构模，则称 \(G\) 的两个表示\textbf{等价}（或相似）。不等价的表示称为\textbf{不等价}的。
\end{definition}

设 \(\varphi : G \to GL(V)\) 与 \(\psi : G \to GL(W)\) 是等价表示（这里 \(V\) 与 \(W\) 必须是同一域 \(F\) 上的向量空间）。设 \(T : V \to W\) 是它们之间的 \(FG\)-模同构。由于 \(T\) 特别地是 \(F\)-模同构，\(T\) 是向量空间同构，故 \(V\) 与 \(W\) 必须有相同的维数。进一步，对所有 \(g \in G\)、\(v \in V\)，由于 \(T\) 是 \(FG\)-模同构有 \(T(g \cdot v) = g \cdot (T(v))\). 按环元素作用的定义这意味着 \(T(\varphi(g)v) = \psi(g)(T(v))\)，即
\[
T \circ \varphi(g) = \psi(g) \circ T \qquad \text{对所有 } g \in G.
\]
特别地，若把 \(V\) 与 \(W\) 作为向量空间等同，则 \(V\) 上 \(G\) 的两个表示 \(\varphi\) 与 \(\psi\) 等价当且仅当存在 \(T \in GL(V)\) 使对所有 \(g \in G\) 有 \(T \circ \varphi(g) \circ T^{-1} = \psi(g)\). 这个 \(T\) 是对所有 \(\varphi(g)\)（\(g \in G\)）的同时基变换。

用矩阵术语，两个表示 \(\varphi\) 与 \(\psi\) 等价当且仅当存在固定的可逆矩阵 \(P\) 使
\[
P\varphi(g)P^{-1} = \psi(g) \qquad \text{对所有 } g \in G.
\]
上面的线性变换 \(T\) 或矩阵 \(P\) 称为\textbf{缠结}（intertwines）表示 \(\varphi\) 与 \(\psi\)（它给出把 \(\varphi\) 变为 \(\psi\) 的「规则」）。

为研究 \(FG\)-模分解为子模（的直和），我们需要一些术语。我们对任意环陈述这些定义，因为下一节将在更一般的情形讨论直和分解。

\begin{definition}
设 \(R\) 是环，\(M\) 是非零 \(R\)-模。

（1）若 \(M\) 的子模只有 0 与 \(M\)，则称 \(M\) 是\textbf{不可约}（或\textbf{单}）的；否则称 \(M\) 是\textbf{可约}的。

（2）若 \(M\) 不能写成 \(M_1 \oplus M_2\)（\(M_1, M_2\) 是非零子模），则称 \(M\) 是\textbf{不可分解}的；否则称 \(M\) 是\textbf{可分解}的。

（3）若 \(M\) 是不可约子模的直和，则称 \(M\) 是\textbf{完全可约}的。

（4）按提供它的 \(FG\)-模是否具有相应性质，称表示是\textbf{不可约}、\textbf{可约}、\textbf{不可分解}、\textbf{可分解}或\textbf{完全可约}的。

（5）若 \(M\) 是完全可约 \(R\)-模，则 \(M\) 的任何直和项称为 \(M\) 的\textbf{成分}（即 \(N\) 是 \(M\) 的成分，若存在 \(M\) 的子模 \(N'\) 使 \(M = N \oplus N'\)）。
\end{definition}

不可约模按定义既不可分解又完全可约。我们很快将给出不可分解但不可约的模的例子。

若 \(R = FG\)，不可约 \(FG\)-模 \(V\) 是没有非平凡真 \(G\)-不变子空间体的非零 \(F\)-向量空间。例如若 \(\dim_FV = 1\)，则 \(V\) 必然不可约（其子空间只有 0 与 \(V\)）。

设 \(V\) 是有限维 \(FG\)-模且 \(V\) 可约。设 \(U\) 是 \(G\)-不变子空间。取 \(U\) 的一组基并把其扩充为 \(V\) 的基。则对每个 \(g \in G\)，\(\varphi(g)\) 关于这组基作用在 \(V\) 上的矩阵具有形式
\[
\varphi(g) = \begin{pmatrix}\varphi_1(g) & \varphi_2(g)\\ 0 & \psi(g)\end{pmatrix},
\]
其中 \(\varphi_1 = \varphi|_U\)（关于 \(U\) 的选定基）而 \(\psi\) 是 \(G\) 在 \(V/U\) 上的表示（且 \(\varphi_2\) 不必是方阵）。故可约表示就是相应的矩阵表示中各矩阵呈分块上三角形式的表示。

再设 \(FG\)-模 \(V\) 可分解，\(V = U \oplus U'\). 取 \(U\) 的一组基与 \(U'\) 的一组基之并作为 \(V\) 的基。在这个基的选取下每个 \(g \in G\) 的矩阵具有形式
\[
\varphi(g) = \begin{pmatrix}\varphi_1(g) & 0\\ 0 & \varphi_2(g)\end{pmatrix}
\]
（即对所有 \(g \in G\) 有 \(\psi(g) = \varphi(g)\)）。故可分解表示就是相应的矩阵表示中各矩阵呈分块对角形式的表示。

\noindent\textbf{例}

（1）如上所述，所有 1 次表示都是不可约、不可分解且完全可约的。特别地，这适用于平凡表示与上面例 5 中描述的表示。

（2）若 \(|G| > 1\)，则 \(G\) 的正则表示可约（增广理想与迹理想是真非零子模）。我们稍后将确定这个表示完全可约的条件以及它如何分解为直和。

（3）对 \(n > 1\)，上面例 10 中描述的 \(FS_n\)-模可约，因为 \(N\) 与 \(I\) 是真非零子模。模 \(N\) 不可约（1 维）；若域 \(F\) 的特征不整除 \(n\)，则 \(I\) 也不可约。

（4）上面例 6 中描述的 \(n \geq 3\) 时二面体群 \(D_{2n} = G\) 的 2 次表示不可约。没有 \(G\)-不变 1 维子空间，因为旋转 \(2\pi/n\) 弧度把 \(\mathbb{R}^2\) 中没有直线映到自身。类似地，例 7 中描述的 \(Q_8\) 的 2 次复表示不可约，因为给定矩阵 \(\varphi(i)\) 恰有两个 1 维特征空间（对应其不同特征值 \(\pm\sqrt{-1}\)），而它们在矩阵 \(\varphi(j)\) 下不变。例 8 中描述的 4 次表示 \(\varphi : Q_8 \to GL_4(\mathbb{R})\) 也可以证明是不可约的（参见习题）。然而我们将看到，若把 \(\varphi\) 看作复表示 \(\varphi : Q_8 \to GL_4(\mathbb{C})\)（即把矩阵的实元素看作复元素），则存在复矩阵 \(P\) 使对所有 \(g \in Q_8\) 有 \(P^{-1}\varphi(g)P\) 是 \(2 \times 2\) 分块矩阵的直和。故域 \(F\) 上的不可约表示在扩域后可能变得可约。

（5）设 \(G = \langle g\rangle\) 是 \(n\) 阶循环群，假设 \(F\) 含有全部 \(n\) 次单位根。如群代数的例子中所述，\(F\langle g\rangle \cong F[x]/(x^n-1)\). 故 \(FG\)-模恰是被 \(x^n-1\) 零化的 \(F[x]\)-模。后者（有限维）模按 Jordan 标准形定理在同构意义下被描述。

若 \(g\) 作用在 \(F\langle g\rangle\)-模 \(V\) 上的极小多项式在 \(F\) 中有互不相同的根，则存在 \(V\) 的一组基使 \(g\)（从而其所有幂）由对角矩阵表示（参见第 12.3 节推论 25）。此时 \(V\) 是完全可约 \(F\langle g\rangle\)-模（作为 1 维 \(\langle g\rangle\)-不变子空间的直和）。一般地，\(g\) 作用在 \(V\) 上的极小多项式整除 \(x^n-1\)，故若 \(x^n-1\) 在 \(F\) 中有互不相同的根，则 \(V\) 是完全可约 \(F\langle g\rangle\)-模。多项式 \(x^n-1\) 在 \(F\) 中有互不相同的根当且仅当 \(F\) 的特征不整除 \(n\). 这给出每个 \(F\langle g\rangle\)-模完全可约的一个充分条件。

若 \(g\) 作用在 \(V\) 上的极小多项式没有互不相同的根（即 \(F\) 的特征整除 \(n\)），则 \(g\) 的 Jordan 标准形必有一个大小大于 1 的初等 Jordan 块。由于每个线性变换有唯一的 Jordan 标准形，\(g\) 不能由对角矩阵表示，即 \(V\) 不完全可约。由第 12.3 节中循环模的结果可知，任何大小大于 1 的 Jordan 块中 \(g\) 的（1 维）特征空间没有 \(\langle g\rangle\)-不变补，即 \(V\) 可约但不完全可约。

具体地，设 \(p\) 是素数，\(F = \mathbb{F}_p\)，\(g\) 的阶为 \(p\). 设 \(V\) 是 \(\mathbb{F}_p\) 上以 \(v, w\) 为基的 2 维空间，并定义 \(g\) 在 \(V\) 上的作用为
\[
g \cdot v = v, \qquad g \cdot w = v+w.
\]
\(V\) 的这个自同态在 \(GL(V)\) 中确实有阶 \(p\)，而 \(g\) 关于这组基的矩阵是初等 Jordan 块
\[
\varphi(g) = \begin{pmatrix}1 & 1\\ 0 & 1\end{pmatrix}.
\]
现在 \(V\) 可约（\(\operatorname{span}\{v\}\) 是 \(\langle g\rangle\)-不变子空间）但 \(V\) 不可分解（上面的 \(2 \times 2\) 初等 Jordan 矩阵不相似于对角矩阵）。

有限群表示论的第一个基本结果表明例 5 如何推广到非循环群。

\begin{theorem}\label{thm:18.1}
（\textbf{Maschke 定理}）设 \(G\) 是有限群，\(F\) 是特征不整除 \(|G|\) 的域。若 \(V\) 是任意 \(FG\)-模，\(U\) 是 \(V\) 的任意子模，则 \(V\) 有子模 \(W\) 使 \(V = U \oplus W\)（即每个子模都是直和项）。
\end{theorem}

\noindent\textbf{注：} 当 \(F\) 的特征为 0 时 Maschke 定理的假设对任何有限群都满足。

\begin{proof}
Maschke 定理证明的想法是构造一个 \(FG\)-模同态
\[
\pi : V \to U,
\]
它是到 \(U\) 上的投射，即满足下述两条性质：

（i）对所有 \(u \in U\) 有 \(\pi(u) = u\)；

（ii）对所有 \(v \in V\) 有 \(\pi(\pi(v)) = \pi(v)\)（即 \(\pi^2 = \pi\)）

（事实上（ii）由（i）与 \(\pi(V) \subseteq U\) 推出）。

先设我们能构造这样的 \(FG\)-模同态，令 \(W = \ker\pi\). 由于 \(\pi\) 是模同态，\(W\) 是子模。我们看出 \(W\) 是 \(U\) 的直和补，如下。若 \(v \in U \cap W\)，则由（i）\(v = \pi(v)\)，而按 \(W\) 的定义 \(\pi(v) = 0\). 这表明 \(U \cap W = 0\). 为证明 \(V = U+W\)，设 \(v\) 是 \(V\) 的任意元素，写 \(v = \pi(v)+(v-\pi(v))\). 按定义 \(\pi(v) \in U\). 由 \(\pi\) 的性质（ii），
\[
\pi(v-\pi(v)) = \pi(v)-\pi(\pi(v)) = \pi(v)-\pi(v) = 0,
\]
即 \(v-\pi(v) \in W\). 这表明 \(V = U+W\)，从而 \(V = U \oplus W\). 因此为建立 Maschke 定理，只需找到这样的 \(FG\)-模投射 \(\pi\).

由于 \(U\) 是子空间，它在 \(V\) 中有向量空间直和补 \(W_0\)（取 \(U\) 的一组基 \(B_1\)，把它扩充为 \(V\) 的基 \(B\)，令 \(W_0\) 是 \(B-B_1\) 的张成）。故作为向量空间有 \(V = U \oplus W_0\)，但 \(W_0\) 不必 \(G\)-稳定（即不必是 \(FG\)-子模）。设 \(\pi_0 : V \to U\) 是这个直和分解相应的向量空间投射，即定义为：对所有 \(u \in U\)、\(w \in W_0\) 有 \(\pi_0(u+w) = u\).

证明的关键想法是对 \(\pi_0\) 在 \(G\) 上「取平均」以构造 \(FG\)-模投射 \(\pi\). 对每个 \(g \in G\) 定义
\[
g\pi_0g^{-1}(v) = g \cdot \pi_0(g^{-1} \cdot v) \qquad \text{对所有 } v \in V
\]
（这里 \(\cdot\) 表示环 \(FG\) 的元素的作用）。由于 \(\pi_0\) 映到 \(U\) 中且在 \(g\) 的作用下稳定，\(g\pi_0g^{-1}\) 把 \(V\) 映到 \(U\). \(g\) 与 \(g^{-1}\) 都作为 \(F\)-线性变换作用，故 \(g\pi_0g^{-1}\) 是线性变换。进一步，若 \(u\) 在 \(G\)-稳定空间 \(U\) 中，则 \(g^{-1}u\) 也在其中，而按 \(\pi_0\) 的定义 \(\pi_0(g^{-1}u) = g^{-1}u\). 由此对每个 \(g \in G\) 与每个 \(u \in U\) 有 \(g\pi_0g^{-1}(u) = u\)（即 \(g\pi_0g^{-1}\) 也是 \(V\) 到 \(U\) 的向量空间投射）。

令 \(n = |G|\)，把 \(n\) 看作 \(F\) 的元素（\(n = 1+\cdots+1\)，\(n\) 个）。由假设 \(n\) 在 \(F\) 中不为零，故在 \(F\) 中有逆。定义
\[
\pi = \frac{1}{n}\sum_{g \in G}g\pi_0g^{-1}.
\]
由于 \(\pi\) 是从 \(V\) 到 \(U\) 的线性变换之和的纯量倍，它也是从 \(V\) 到 \(U\) 的线性变换。进一步，定义 \(\pi\) 的和中每一项都限制为子空间 \(U\) 上的恒等映射，故 \(\pi|_U\) 是 \(n\) 个恒等映射之和的 \(1/n\) 倍。这些观察证明
\[
\pi : V \to U \text{ 是线性变换；} \quad \pi(u) = u \text{ 对所有 } u \in U； \quad \pi^2(v) = \pi(v) \text{ 对所有 } v \in V.
\]
还需证明 \(\pi\) 是 \(FG\)-模同态（即 \(FG\)-线性的）。只需证明对所有 \(h \in G\)、\(v \in V\) 有 \(\pi(hv) = h\pi(v)\). 此时
\[
\pi(hv) = \frac{1}{n}\sum_{g \in G}g\pi_0(g^{-1}hv) = \frac{1}{n}\sum_{k \in G}hk\pi_0(k^{-1}v) = h\pi(v)
\]
（当 \(g\) 取遍 \(G\) 的所有元素时 \(k = h^{-1}g\) 也如此，且模元素 \(h\) 可以由模中的分配律提到求和之外）。这建立了 \(FG\)-模投射 \(\pi\) 的存在性，从而完成证明。
\end{proof}

Maschke 定理的应用将针对有限生成 \(FG\)-模。然而与 \(F[x]\)-模的情形不同，有限生成 \(FG\)-模自动是有限维向量空间（差别在于 \(FG\) 本身有限维，而 \(F[x]\) 不是）。设 \(V\) 是 \(FG\)-模。若 \(V\) 是 \(F\) 上有限维向量空间，则 \(V\) 作为 \(FG\)-模当然有限生成（任何 \(F\)-基给出 \(FG\) 上的一组生成元）。反之，若 \(V\) 作为 \(FG\)-模有限生成，比如由 \(v_1, \ldots, v_k\) 生成，则容易看出 \(V\) 作为向量空间由有限集 \(\{g \cdot v_i \mid g \in G,\ 1 \leq i \leq k\}\) 张成。故

\noindent\textbf{\(FG\)-模有限生成当且仅当它有限维。}

\begin{corollary}\label{cor:18.2}
若 \(G\) 是有限群，\(F\) 是特征不整除 \(|G|\) 的域，则每个有限生成 \(FG\)-模都完全可约（等价地，\(G\) 的每个有限次 \(F\)-表示都完全可约）。
\end{corollary}

\begin{proof}
设 \(V\) 是有限生成 \(FG\)-模。如上所述 \(V\) 在 \(F\) 上有限维，故我们可以对其维数归纳。若 \(V\) 不可约，它完全可约，结论成立。否则设 \(V\) 有真非零 \(FG\)-子模 \(U\). 由 Maschke 定理 \(U\) 有 \(FG\)-子模补 \(W\)，即 \(V = U \oplus W\). 由归纳假设 \(U\) 与 \(W\) 都是不可约子模的直和，故 \(V\) 也是。这完成了归纳。
\end{proof}

\begin{corollary}\label{cor:18.3}
设 \(G\) 是有限群，\(F\) 是特征不整除 \(|G|\) 的域，\(\varphi : G \to GL(V)\) 是 \(G\) 的有限次表示。则存在 \(V\) 的一组基，使对每个 \(g \in G\)，\(\varphi(g)\) 关于这组基的矩阵是分块对角的，且各分块 \(\varphi_i(g)\)（\(1 \leq i \leq m\)）是 \(G\) 的不可约矩阵表示。
\end{corollary}

\begin{proof}
由推论 2 我们可以写 \(V = U_1 \oplus U_2 \oplus \cdots \oplus U_m\)，其中 \(U_i\) 是 \(V\) 的不可约 \(FG\)-子模。设 \(B_i\) 是 \(U_i\) 的基，\(B\) 是各 \(B_i\) 的并。对每个 \(g \in G\)，\(\varphi(g)\) 关于基 \(B\) 的矩阵具有推论中的形式，其中 \(\varphi_i(g)\) 是 \(\varphi(g)|_{U_i}\) 关于基 \(B_i\) 的矩阵。
\end{proof}

Maschke 定理的逆也成立。即若 \(F\) 的特征整除 \(|G|\)，则 \(G\) 拥有不完全可约的（有限生成）\(FG\)-模。特别地，正则表示（即模 \(FG\) 本身）不完全可约。

在第 18.2 节我们将讨论 \(FG\)-模分解为不可约子模的直和时各成分的唯一性问题。

\subsection*{习题}

设 \(F\) 是域，\(G\) 是有限群，\(n \in \mathbb{Z}^+\).

\begin{enumerate}
\item 证明若 \(\varphi : G \to GL(V)\) 是任意表示，则 \(\varphi\) 给出 \(G/\ker\varphi\) 的忠实表示。

\item 设 \(\varphi : G \to GL_n(F)\) 是矩阵表示。证明映射 \(g \mapsto \det(\varphi(g))\) 是 1 次表示。

\item 证明 \(G\) 的 1 次表示与阿贝尔群 \(G/G'\) 的 1 次表示一一对应（其中 \(G'\) 是 \(G\) 的换位子群）。

\item 设 \(V\) 是 \(FG\)-模（可能无限维，\(G\) 是有限群）。证明对每个 \(v \in V\) 存在包含 \(v\) 的维数 \(\leq |G|\) 的 \(FG\)-子模。

\item 证明若 \(|G| > 1\)，则每个不可约 \(FG\)-模的维数都小于 \(|G|\).

\item 对下列在第二组例子中描述的每个表示，写出所有 \(g \in G\) 的矩阵 \(\varphi(g)\)：

（a）例 3 中描述的 \(S_3\) 的表示（取该例中 \(n = 3\)）；

（b）例 6 中描述的 \(D_8\) 的表示（即取该例中 \(n = 4\) 并写出所有群元素的所有正弦、余弦值）；

（c）例 7 中描述的 \(Q_8\) 的表示；

（d）例 8 中描述的 \(Q_8\) 的表示。

\item 设 \(V\) 是第二组例子中例 3 描述的 \(S_4\) 的 4 维置换模。设 \(\pi : D_8 \to S_4\) 是由 \(D_8\) 左乘作用在其子群 \(\langle s\rangle\) 的左陪集之集上所得的 \(D_8\) 的置换表示。如例 4 中描述的那样通过 \(\pi\) 把 \(V\) 变成 \(FD_8\)-模，并写出这个表示给出的 \(r\) 与 \(s\) 关于基 \(e_1, \ldots, e_4\) 的 \(4 \times 4\) 矩阵。

\item 设 \(V\) 是第二组例子中例 3 与例 10 描述的 \(FS_n\)-模。

（a）证明若 \(v\) 是 \(V\) 中满足对所有 \(\sigma \in S_n\) 有 \(\sigma \cdot v = v\) 的任意元素，则 \(v\) 是 \(e_1+e_2+\cdots+e_n\) 的 \(F\)-倍。

（b）证明若 \(n \geq 3\)，则 \(V\) 有唯一的 1 维子模，即由 \(e_1+e_2+\cdots+e_n\) 的所有 \(F\)-倍构成的子模 \(N\).

\item 证明第二组例子中例 8 描述的 \(Q_8\) 在 \(\mathbb{H}\) 上的 4 维表示不可约。〔证明任何 \(Q_8\)-稳定子空间都是左理想。〕

\item 证明 \(GL_2(\mathbb{R})\) 没有同构于 \(Q_8\) 的子群。〔可以用 \(Q_8\) 的生成元与关系直接计算。把其中一个生成元取为有理标准形可简化这些计算。〕

\item 设 \(\varphi : S_n \to GL_n(F)\) 是第二组例子中例 3 描述的置换模给出的矩阵表示（矩阵关于基 \(e_1, \ldots, e_n\) 计算）。证明对所有 \(\sigma \in S_n\) 有 \(\det\varphi(\sigma) = \varepsilon(\sigma)\)，其中 \(\varepsilon(\sigma)\) 是置换 \(\sigma\) 的符号。〔在对换上验证。〕

\item 假设 \(F\) 的特征不为 2. 设 \(H\) 是使 \(T\) 的每一行与每一列恰有一个非零元素、其他位置为零且非零元素为 \(\pm1\) 的 \(T \in M_n(F)\) 之集。证明 \(H\) 是 \(GL_n(F)\) 的子群且 \(H \cong E_{2^n} \rtimes S_n\)（半直积），其中 \(E_{2^n}\) 是 \(2^n\) 阶初等阿贝尔群。

\noindent 下面几个习题探索称为 Schur 引理的重要结果及其一些推论。

\item 设 \(R\) 是环，\(M\) 与 \(N\) 是单（即不可约）\(R\)-模。

（a）证明每个非零 \(R\)-模同态 \(M \to N\) 都是同构。〔考虑其核与像。〕

（b）证明 \textbf{Schur 引理}：若 \(M\) 是单 \(R\)-模，则 \(\operatorname{Hom}_R(M, M)\) 是除环（回忆 \(\operatorname{Hom}_R(M,M)\) 是从 \(M\) 到 \(M\) 的所有 \(R\)-模同态之环，其中乘法是复合）。

\item 设 \(\varphi : G \to GL(V)\) 是 \(G\) 的表示。\(\varphi\) 的\textbf{中心化子}定义为满足对所有 \(g \in G\) 有 \(A\varphi(g) = \varphi(g)A\) 的所有 \(V\) 到自身的线性变换 \(A\) 之集（即与所有 \(\varphi(g)\) 交换的 \(V\) 的线性变换）。

（a）证明 \(V\) 到自身的线性变换 \(A\) 在 \(\varphi\) 的中心化子中当且仅当它是 \(V\) 到自身的 \(FG\)-模同态（故 \(\varphi\) 的中心化子与环 \(\operatorname{Hom}_{FG}(V,V)\) 相同）。

（b）证明若 \(z\) 在 \(G\) 的中心中，则 \(\varphi(z)\) 在 \(\varphi\) 的中心化子中。

（c）假设 \(\varphi\) 是不可约表示（故 \(V\) 是单 \(FG\)-模）。证明若 \(H\) 是 \(GL(V)\) 的满足「\(A \in H\) 时对所有 \(g\) 有 \(A\varphi(g) = \varphi(g)A\)」的有限阿贝尔子群，则 \(H\) 循环（换句话说，环 \(\operatorname{Hom}_{FG}(V,V)\) 的单位乘法群的任何有限阿贝尔子群都循环）。〔由上一习题 \(\operatorname{Hom}_{FG}(V,V)\) 是除环，故这化归为证明除环中非零元素的乘法群的有限阿贝尔子群循环。证明任何除环的阿贝尔子群生成的除子环是域，并利用第 9.5 节命题 18。〕

（d）证明若 \(\varphi\) 是忠实的不可约表示，则 \(G\) 的中心循环。

（e）由（d）推出若 \(G\) 阿贝尔且 \(\varphi\) 是任意不可约表示，则 \(G/\ker\varphi\) 循环。

\item 给出有限循环群的所有 1 维复表示；并确定哪些不等价。

\item 给出有限阿贝尔群的所有 1 维复表示。由此证明有限阿贝尔群的不等价 1 次复表示的个数等于群的阶。〔先把阿贝尔群分解为循环群的直积，然后用上一习题。〕

\item 证明阿贝尔群的复表示的下述 Schur 引理变形：若 \(G\) 阿贝尔，则 \(G\) 的任何不可约复表示 \(\varphi\) 的次数为 1，且 \(G/\ker\varphi\) 循环。〔这可以不引用习题 14，而利用「对任意 \(g \in G\)，\(\varphi(g)\) 的特征空间是 \(G\)-稳定的」这一观察来完成。你关于 \(\varphi\) 次数为 1 的证明也应当对无限阿贝尔群成立。〕

\item 证明复表示的下述一般形式的 Schur 引理：若 \(\varphi : G \to GL_n(\mathbb{C})\) 是不可约矩阵表示，\(A\) 是与所有 \(\varphi(g)\)（\(g \in G\)）交换的 \(n \times n\) 矩阵，则 \(A\) 是纯量矩阵。由此证明若 \(\varphi\) 是忠实的、不可约的复表示，则 \(G\) 的中心循环，且对 \(G\) 的中心中所有元素 \(z\)，\(\varphi(z)\) 是纯量矩阵。〔如上一习题，\(A\) 的特征空间是 \(G\)-稳定的。〕

\item 证明若 \(G\) 是阿贝尔群，则 \(G\) 的任何有限维复表示都等价于到对角矩阵的表示（即任何有限个两两交换的复矩阵可以同时对角化）。〔这可以不引用 Maschke 定理，而通过考察特征空间来完成。〕

\item 证明任何有限群 \(G\) 的 1 次复表示的个数等于 \(|G:G'|\)，其中 \(G'\) 是 \(G\) 的换位子群。〔利用习题 3 与 16。〕

\item 设 \(G\) 是非循环阿贝尔群，通过共轭作用在初等阿贝尔 \(p\)-群 \(V\) 上，其中 \(p\) 是不整除 \(G\) 的阶的素数。

（a）证明若 \(W\) 是 \(V\) 的不可约 \(\mathbb{F}_pG\)-子模，则存在某个非单位元 \(g \in G\) 使 \(W \subseteq C_V(g)\)（这里 \(C_V(g)\) 是共轭下被 \(g\) 固定的 \(V\) 的元素之集）。

（b）证明当 \(g\) 取遍 \(G\) 的所有非单位元时，\(V\) 由各子群 \(C_V(g)\) 生成。

\item 设 \(p\) 是素数，\(P\) 是 \(p\)-群，\(F\) 是特征为 \(p\) 的域。证明 \(P\) 在 \(F\) 上唯一的不可约表示是平凡表示。〔先对 \(p\) 阶群证明：利用 \(F\) 含有全部 \(p\) 次单位根（即 1 本身）。若 \(P\) 不是 \(p\) 阶，设 \(z\) 是 \(P\) 的中心中 \(p\) 阶元素，证明 \(z\) 在不可约表示的核中，并对 \(P/\langle z\rangle\) 归纳。〕

\item 设 \(p\) 是素数，\(P\) 是非平凡 \(p\)-群，\(F\) 是特征为 \(p\) 的域。证明正则表示不完全可约。〔利用上一习题。〕

\item 设 \(p\) 是素数，\(P\) 是非平凡 \(p\)-群，\(F\) 是特征为 \(p\) 的域。证明正则表示不可分解。
\end{enumerate}

\section{Wedderburn 定理及其若干推论}

本节给出一个著名的分类定理，它由 Wedderburn 证明，特别地描述了当 \(F\) 的特征不整除群 \(G\) 的阶时群代数 \(FG\) 的结构。由这个分类定理我们将导出各种推论，其中包括：对每个有限群 \(G\)，只有有限多个互不同构的不可约 \(FG\)-模。这一结果连同 Maschke 定理，在某种意义上完成了有限群表示论在这类域上的 Hölder 纲领。本书余下部分致力于发展确定和处理不可约表示的技巧，以及把这些知识应用于获得群论信息。

\begin{theorem}\label{thm:18.4}
（\textbf{Wedderburn 定理}）设 \(R\) 是非零环（不必交换）。则下列条件等价：

（1）每个 \(R\)-模都是投射的；

（2）每个 \(R\)-模都是内射的；

（3）每个 \(R\)-模都是完全可约的；

（4）把 \(R\) 看作左 \(R\)-模，它是直和
\[
R = L_1 \oplus L_2 \oplus \cdots \oplus L_n,
\]
其中每个 \(L_i\) 是单模（即单左理想），且对某个 \(e_i \in R\) 有 \(L_i = Re_i\)，满足

（i）当 \(i \neq j\) 时 \(e_ie_j = 0\)；

（ii）对所有 \(i\) 有 \(e_i^2 = e_i\)；

（iii）\(\sum_{i=1}^{n}e_i = 1\)；

（5）作为环，\(R\) 同构于除环上矩阵环的直积，即
\[
R = R_1 \times R_2 \times \cdots \times R_r,
\]
其中 \(R_j\) 是 \(R\) 的双边理想，且 \(R_j\) 同构于元素取自除环 \(\Delta_j\) 的所有 \(n_j \times n_j\) 矩阵之环，\(j = 1, 2, \ldots, r\). 整数 \(r\)、整数 \(n_j\) 以及除环 \(\Delta_j\)（在同构意义下）都由 \(R\) 唯一确定。
\end{theorem}

\begin{proof}
Wedderburn 定理的证明在习题 1 至 10 中给出概要。
\end{proof}

\begin{definition}
称满足定理 4 中任一条（等价）性质的环 \(R\) 为\textbf{带极小条件的半单环}。
\end{definition}

满足定理 4 中任一等价条件的环 \(R\) 也满足左理想上的\textbf{极小条件}（或\textbf{降链条件}，D.C.C.）：

若 \(I_1 \supseteq I_2 \supseteq \cdots\) 是 \(R\) 的左理想的降链，则存在 \(N \in \mathbb{Z}^+\)，使对所有 \(k \geq N\) 有 \(I_k = I_N\).

（这也解释了上述定义中术语的用法。）我们处理的环都具有这个极小条件。例如群代数总有此性质，因为在任何严格降链中，理想的向量空间维数（它们是 \(FG\) 的 \(F\)-子空间）严格递减，故严格降链的长度至多等于 \(\dim_F FG\)（\(= |G|\)）。因此我们将用术语「半单」表示「带极小条件的半单」。定理 4 结论（5）中的环 \(R_j\) 称为 \(R\) 的 \textbf{Wedderburn 分量}，而 \(R\) 的直积分解称为它的 \textbf{Wedderburn 分解}。注意交换环的 Wedderburn 定理是第 16.1 节 Artin 环分类的推论。带极小条件的交换半单环是 Jacobson 根为零的 Artin 环，故是域的直积（这些域是它的 Wedderburn 分量）。

应当注意，条件（5）是一个双边条件，它描述了 \(R\) 的整体结构（这个环的直积中的环运算是逐分量加法与乘法）。特别地，它蕴含半单环在右理想上也满足极小条件。思考结论（5）中直积 \(R_1 \times \cdots \times R_r\) 的元素的一个有用方式是把它们看作形如
\[
\begin{pmatrix}
A_1 & & \\
& \ddots & \\
& & A_r
\end{pmatrix}
\]
的 \(n \times n\)（分块对角）矩阵，其中 \(A_i\) 是元素取自 \(\Delta_i\) 的任意 \(n_i \times n_i\) 矩阵（这里 \(n = \sum_{i=1}^{r}n_i\)）。

回忆第 10.5 节，若只要 \(Q\) 是任意 \(R\)-模 \(M\) 的子模，\(M\) 就有子模 \(N\) 使 \(M = Q \oplus N\)，则称 \(R\)-模 \(Q\) 是内射的。故 Maschke 定理蕴含：

\begin{corollary}\label{cor:18.5}
若 \(G\) 是有限群，\(F\) 是特征不整除 \(|G|\) 的域，则群代数 \(FG\) 是半单环。
\end{corollary}

在获得关于不变量 \(n_j\)、\(r\)、\(\Delta_j\) 等如何与某些域 \(F\) 的群环 \(FG\) 中的不变量相关的更精确信息之前，我们先研究矩阵环（即 Wedderburn 定理结论（4）与（5）中描述的环）的结构。我们引入一些在环论中广泛使用的术语。

回忆环 \(R\) 的\textbf{中心}是与 \(R\) 中所有元素交换的元素构成的子环；记作 \(Z(R)\)（若环含 1，则中心包含 1）。

\begin{definition}
（1）环 \(R\) 中的非零元素 \(e\) 若满足 \(e^2 = e\)，则称为\textbf{幂等元}。

（2）若 \(e_1e_2 = e_2e_1 = 0\)，则称幂等元 \(e_1\) 与 \(e_2\) \textbf{正交}。

（3）若幂等元 \(e\) 不能写成两个（交换的）正交幂等元之和，则称它是\textbf{本原的}。

（4）若 \(e \in Z(R)\)，且 \(e\) 不能写成环 \(Z(R)\) 中两个正交幂等元之和，则称幂等元 \(e\) 为\textbf{本原中心幂等元}。
\end{definition}

命题 6 描述矩阵环的理想结构，命题 8 把这些结果推广到矩阵环的直积。

\begin{proposition}\label{pro:18.6}
设 \(\Delta\) 是除环，\(n \in \mathbb{Z}^+\)，\(R\) 是元素取自 \(\Delta\) 的所有 \(n \times n\) 矩阵之环，\(I\) 是单位矩阵（即 \(R\) 的 1）。

（1）\(R\) 的双边理想只有 0 与 \(R\)。

（2）\(R\) 的中心由纯量矩阵 \(aI\) 构成，其中 \(a\) 在 \(\Delta\) 的中心中：\(Z(R) = \{aI \mid a \in Z(\Delta)\}\)，且这是与 \(Z(\Delta)\) 同构的域。特别地，若 \(\Delta\) 是域，则 \(R\) 的中心是全体纯量矩阵构成的子环。\(R\) 中唯一的中心幂等元是 \(I\)（特别地，\(I\) 是本原的）。

（3）设 \(e_i\) 是在位置 \(i,i\) 处为 1、其余位置为 0 的矩阵。则 \(e_1, \ldots, e_n\) 是正交本原幂等元，且 \(\sum_{i=1}^{n}e_i = I\).

（4）\(L_i = Re_i\) 是由第 \(i\) 列元素任意、其余列全为零的矩阵构成的左理想。\(L_i\) 是单左 \(R\)-模。每个单左 \(R\)-模都同构于 \(L_1\)（特别地，所有 \(L_i\) 作为 \(R\)-模同构），且作为左 \(R\)-模有 \(R = L_1 \oplus \cdots \oplus L_n\).
\end{proposition}

在证明这个命题之前，先给出下述结果会很有用。

\begin{lemma}\label{lem:18.7}
设 \(R\) 是任意非零环。

（1）若 \(M\) 与 \(N\) 是单 \(R\)-模，\(\varphi : M \to N\) 是非零 \(R\)-模同态，则 \(\varphi\) 是同构。

（2）（\textbf{Schur 引理}）若 \(M\) 是单 \(R\)-模，则 \(\operatorname{Hom}_R(M, M)\) 是除环。
\end{lemma}

\begin{proof}
为证明（1），注意由于 \(\varphi\) 非零，\(\ker\varphi\) 是 \(M\) 的真子模。由 \(M\) 的单性有 \(\ker\varphi = 0\). 类似地，\(\varphi\) 的像是单模 \(N\) 的非零子模，故 \(\varphi(M) = N\). 这证明 \(\varphi\) 是双射，从而（1）成立。

由（1），环 \(\operatorname{Hom}_R(M, M)\) 的每个非零元素都是同构，故有逆。这给出（2）。
\end{proof}

\begin{proof}[命题 6 的证明]
设 \(A\) 是 \(R\) 中 \(i,j\) 位置为 \(a_{ij}\) 的任意矩阵。设 \(E_{ij}\) 是在位置 \(i,j\) 处为 1、其余位置为 0 的矩阵。下列直接计算留作习题：

（i）\(E_{ij}A\) 是第 \(i\) 行等于 \(A\) 的第 \(j\) 行、其余行全为零的矩阵；

（ii）\(AE_{ij}\) 是第 \(j\) 列等于 \(A\) 的第 \(i\) 列、其余列全为零的矩阵；

（iii）\(E_{pq}AE_{rs}\) 是 \(p,s\) 位置为 \(a_{qr}\)、其余位置全为零的矩阵。

为证明（1），设 \(J\) 是 \(R\) 的任意非零双边理想，\(A\) 是 \(J\) 中在位置 \(q,r\) 有非零元素的元素。给定任意 \(p, s \in \{1, \ldots, n\}\)，由（iii）得到
\[
E_{ps} = \frac{1}{a_{qr}}E_{pq}AE_{rs} \in J.
\]
由于 \(\{E_{ps} \mid 1 \leq p \leq n,\ 1 \leq s \leq n\}\) 的 \(\Delta\)-线性组合给出整个 \(R\)，推出 \(J = R\). 这证明（1）。

为证明（2），设 \(A \in Z(R)\). 则对所有 \(i,j\) 有 \(E_{ij}A = AE_{ij}\). 由上文的（i）与（ii）立刻推出 \(A\) 的所有非对角元素为零，且所有对角元素相等。故 \(A = aI\)，其中 \(a \in \Delta\). 进一步，\(A\) 还必须与所有纯量矩阵之集 \(\{\beta I \mid \beta \in \Delta\}\) 交换，即 \(a\) 必须与 \(\Delta\) 的所有元素交换。最后，由于 \(Z(R)\) 是域，立刻可知它含唯一的幂等元（即 \(I\)）。这建立了（2）的所有部分。

对于（3），显然 \(e_1, \ldots, e_n\) 是正交幂等元，其和为 \(I\). 我们把它们本原的证明推迟到建立（4）之后。

下面证明（4）。由上面的（ii）推出 \(Re_i = RE_{ii}\) 是所有矩阵之集，其中第 \(i\) 列元素任意、其余列全为零。进一步，若 \(A\) 是 \(Re_i\) 的任意非零元素，则显然 \(RA \subseteq Re_i\). 反向包含也成立，因为若 \(a_{pi}\) 是 \(A\) 的非零元素，则由上面的（i）
\[
e_i = E_{ii} = \frac{1}{a_{pi}}E_{ip}A \in RA.
\]
这证明对 \(Re_i\) 中任意非零元素 \(A\) 有 \(Re_i = RA\)，故 \(Re_i\) 必是单 \(R\)-模。

设 \(M\) 是任意单 \(R\)-模。由于对所有 \(m \in M\) 有 \(1m = m\)，且 \(1 = \sum_{i=1}^{n}e_i\)，存在某个 \(i\) 与某个 \(m \in M\) 使 \(e_im \neq 0\). 对这个 \(i\) 与这个 \(m\)，映射 \(re_i \mapsto re_im\) 是从单 \(R\)-模 \(Re_i\) 到单模 \(M\) 的非零 \(R\)-模同态。由引理 7（1）它是同构。由（ii），映射 \(r \mapsto rE_{i1}\) 给出 \(Re_i \cong Re_1\). 最后，每个矩阵都是它的各列之直和，故 \(R = L_1 \oplus \cdots \oplus L_n\). 这完成了（4）的证明。

剩下要证明（3）中的幂等元是本原的。若 \(e_i = a+b\)，其中 \(a,b\) 是正交幂等元，则我们将看到
\[
L_i = Re_i = Ra \oplus Rb.
\]
这将与 \(L_i\) 是单 \(R\)-模这一事实矛盾。为建立上述直和，先注意由于 \(ab = ba = 0\)，有 \(ae_i = a \in Re_i\) 且 \(be_i = b \in Re_i\). 对所有 \(r \in R\) 有 \(re_i = ra+rb\)，故 \(Re_i = Ra+Rb\). 此外 \(Ra \cap Rb = 0\)，因为若对某个 \(r,s \in R\) 有 \(ra = sb\)，则 \(ra = raa = sba = 0\)（回忆 \(a = a^2\) 且 \(ba = 0\)）。这完成了证明的所有部分。
\end{proof}

\begin{proposition}\label{pro:18.8}
设 \(R = R_1 \times R_2 \times \cdots \times R_r\)，其中 \(R_i\) 是除环 \(\Delta_i\) 上 \(n_i \times n_i\) 矩阵之环，\(i = 1, 2, \ldots, r\).

（1）把 \(R_i\) 与直积的第 \(i\) 个分量等同。设 \(z_i\) 是在第 \(i\) 个位置为 \(R_i\) 的单位元、其余位置为零的 \(r\) 元组。则 \(R_i = z_iR\)，且对任意 \(a \in R_i\) 有 \(z_ia = a\)，而对所有 \(j \neq i\) 有 \(z_ja = 0\). 元素 \(z_1, \ldots, z_r\) 就是 \(R\) 的全部本原中心幂等元。它们两两正交，且 \(\sum_{i=1}^{r}z_i = 1\).

（2）设 \(N\) 是任意左 \(R\)-模，令 \(z_iN = \{z_ix \mid x \in N\}\)，\(1 \leq i \leq r\). 则 \(z_iN\) 是 \(N\) 的左 \(R\)-子模，每个 \(z_iN\) 是 \(R_i\)-模，且对所有 \(j \neq i\)，\(R_j\) 在其上平凡作用，并且
\[
N = z_1N \oplus z_2N \oplus \cdots \oplus z_rN.
\]

（3）单 \(R\)-模就是下述意义下的单 \(R_i\)-模：对所有 \(j \neq i\)，\(R_j\) 在其上平凡作用。设 \(M_i\) 是唯一的单 \(R_i\)-模（参看命题 6）。通过让所有 \(j \neq i\) 的 \(R_j\) 平凡作用，可把 \(M_i\) 看作 \(R\)-模。则 \(M_1, \ldots, M_r\) 是两两不同构的单 \(R\)-模，且任何单 \(R\)-模都同构于 \(M_1, \ldots, M_r\) 之一。确切地说，\(R\)-模 \(M_i\) 同构于 \(R\) 的单左理想 \((0, \ldots, 0, L^{(i)}, 0, \ldots, 0)\)，即第 \(i\) 个分量 \(L^{(i)}\) 由第一列元素任意、其余列全为零的矩阵构成的那些元素之集。

（4）对任意 \(R\)-模 \(N\)，\(R\)-子模 \(z_iN\) 是单 \(R\)-模的直和，其中每一个都同构于（3）中的模 \(M_i\). 特别地，若 \(M\) 是单 \(R\)-模，则存在唯一的 \(i\) 使 \(z_iM = M\)，且对这个指标 \(i\)，对所有 \(j \neq i\) 有 \(M \cong M_i\) 以及 \(z_jM = 0\).

（5）若每个 \(\Delta_i\) 都等于域 \(F\)，则 \(R\) 是 \(F\) 上维数为 \(\sum_{i=1}^{r}n_i^2\) 的向量空间，且 \(\dim_FZ(R) = r\).
\end{proposition}

\begin{proof}[命题 8 的证明]
在（1）中，由于环的直积中的乘法是逐分量的，显然 \(z_i\) 与 \(R\) 的元素 \((a_1, \ldots, a_r)\) 之积是在第 \(i\) 个位置为 \(a_i\)、其余位置为零的 \(r\) 元组。故 \(R_i = z_iR\)，\(z_i\) 是 \(R_i\) 的单位元，且对任意 \(j \neq i\)，若 \(a \in R_j\) 则 \(z_ia = 0\). 同样显然 \(z_1, \ldots, z_r\) 是两两正交的中心幂等元，其和为 \(R\) 的单位元。按定义，\(R\) 的中心幂等元是 \(Z(R) = F_1 \times F_2 \times \cdots \times F_r\) 中的幂等元，其中 \(F_i\) 是 \(R_i\) 的中心。由命题 6，\(F_i\) 是域 \(Z(\Delta_i)\). 若 \(w = (w_1, \ldots, w_r)\) 是任意中心幂等元，则对所有 \(i\) 有 \(w_i \in F_i\)，且由 \(w^2 = w\) 得 \(w_i^2 = w_i\)；而域中方程 \(x^2 = x\) 只有解 0 与 1，故 \(R\) 的唯一的中心幂等元是各分量取 0 或 1 的 \(r\) 元组。因此 \(z_1, \ldots, z_r\) 是本原中心幂等元，且由于每个中心幂等元都是它们之和，它们是 \(R\) 的全部本原中心幂等元。这证明（1）。

为证明（2），设 \(N\) 是任意左 \(R\)-模。先注意对任意 \(z \in Z(R)\)，集合 \(\{zx \mid x \in N\}\) 是 \(N\) 的 \(R\)-子模。特别地，\(z_iN\) 是 \(R\)-子模。设 \(z_ix \in z_iN\)，\(a \in R_j\)（其中 \(j \neq i\)）。由（1）有 \(a = az_j\)，故 \(az_ix = (az_j)(z_ix) = az_jz_ix = 0\)，因为 \(z_jz_i = 0\). 因此对所有 \(j \neq i\)，\(R_j\) 在 \(R\)-子模 \(z_iN\) 上平凡作用。对每个 \(x \in N\)，由（1）有 \(x = 1x = z_1x + \cdots + z_rx\)，故 \(N = z_1N + \cdots + z_rN\). 最后，这个和是直和，因为例如若 \(x \in z_1N \cap (z_2N + \cdots + z_rN)\)，则 \(x = z_1x\)，而 \(z_1\) 与 \(z_2N + \cdots + z_rN\) 的任意元素之积为零。这证明（2）。

在（3）中，先注意当对所有 \(j \neq i\) 定义 \(R_j\) 在 \(R_i\)-模 \(M\) 上平凡作用时，\(M\) 成为 \(R\)-模。对这种模 \(M\)，\(R\)-子模与 \(R_i\)-子模相同。故对每个 \(i\)，由于 \(M_i\) 是单 \(R_i\)-模，它是单 \(R\)-模。

其次，设 \(M\) 是单 \(R\)-模。由（2），\(M = z_1M \oplus \cdots \oplus z_rM\). 由于 \(M\) 没有非平凡真 \(R\)-子模，必存在唯一的 \(i\) 使 \(M = z_iM\)，且对所有 \(j \neq i\) 有 \(z_jM = 0\). 故单 \(R\)-模 \(M\) 被所有 \(j \neq i\) 的 \(R_j\) 零化。这意味着 \(M\) 的 \(R\)-子模与 \(M\) 的 \(R_i\)-子模相同，故 \(M\) 是单 \(R_i\)-模。由命题 6，作为 \(R_i\)-模 \(M \cong M_i\). 由于对所有 \(j \neq i\)，\(R_j\) 在 \(M\) 与 \(M_i\) 上都平凡作用，推出该 \(R_i\)-模同构也可看作 \(R\)-模同构。

假设 \(i \neq j\)，且 \(\varphi : M_i \to M_j\) 是 \(R\)-模同构。若 \(s_i \in M_i\)，则 \(s_i = z_is_i\)，故
\[
\varphi(s_i) = \varphi(z_is_i) = z_i\varphi(s_i) = 0,
\]
因为 \(\varphi(s_i) \in M_j\) 而 \(z_i\) 在 \(M_j\) 上平凡作用。这与 \(\varphi\) 是同构矛盾，从而证明 \(M_1, \ldots, M_r\) 是两两不同构的单 \(R\)-模。

最后，（3）中描述的 \(R\) 的左理想被所有 \(j \neq i\) 的 \(R_j\) 平凡作用；且由命题 6，它在同构意义下是唯一的单 \(R_i\)-模。故这个左理想是单 \(R\)-模，且同构于 \(M_i\). 这证明（3）。

对于（4），我们已证明若 \(M\) 是任意单 \(R\)-模，则存在唯一的 \(i\) 使 \(z_iM = M\)，且对所有 \(j \neq i\) 有 \(z_jM = 0\). 进一步，我们已证明对这个指标 \(i\)，单 \(R\)-模 \(M\) 同构于 \(M_i\). 现在设 \(N\) 是任意 \(R\)-模。则 \(z_iN\) 是 \(R_i\)-模，且被所有 \(j \neq i\) 的 \(R_j\) 平凡作用。由 Wedderburn 定理，\(z_iN\) 是单 \(R\)-模的直和。由于这些单直和项中每一个都被所有 \(j \neq i\) 的 \(R_j\) 平凡作用，每一个都同构于 \(M_i\). 这证明（4）。

在（5）中，若每个 \(\Delta_i\) 都等于域 \(F\)，则作为 \(F\)-向量空间，每个矩阵环 \(M_{n_i}(F)\) 在 \(F\) 上维数为 \(n_i^2\)，故 \(R\) 在 \(F\) 上维数为 \(\sum_{i=1}^{r}n_i^2\). 进一步，每个 \(M_{n_i}(F)\) 的中心是 1 维的（因为由命题 6（2）它同构于 \(F\)），故 \(Z(R)\) 在 \(F\) 上维数为 \(r\). 这完成了命题的证明。
\end{proof}

现在把 Wedderburn 定理（以及上述环论计算）应用于群代数 \(FG\). 首先，为了应用 Wedderburn 定理，我们需要 \(F\) 的特征不整除 \(|G|\). 事实上，由于后继各节将处理数值数据，方便起见我们让 \(F\) 的特征等于 0. 其次，若我们迫使出现在 \(FG\) 的 Wedderburn 分解中的所有除环都等于域 \(F\)，事情会简化——我们将证明施加 \(F\) 代数闭这一条件足以保证这一点。为简化记号，本书余下大部分取 \(F = \mathbb{C}\). 读者容易验证，任何特征为 0 的代数闭域（例如全体代数数构成的域）都可在全文中代替 \(\mathbb{C}\).

由推论 5，环 \(\mathbb{C}G\) 是半单的，故由 Wedderburn 定理
\[
\mathbb{C}G \cong R_1 \times R_2 \times \cdots \times R_r,
\]
其中 \(R_i\) 是除环 \(\Delta_i\) 上 \(n_i \times n_i\) 矩阵之环。把这个直积的元素看作 \(n \times n\) 分块矩阵（\(n = \sum_{i=1}^{r}n_i\)），其中第 \(i\) 块元素取自 \(\Delta_i\)，则域 \(\mathbb{C}\) 在这个直积中作为纯量矩阵出现，并含于 \(\mathbb{C}G\) 的中心，从而作为纯量也含于每个 \(\Delta_i\) 的中心 \(Z(\Delta_i)\). 注意每个 \(\Delta_i\) 都是 \(\mathbb{C}\) 上的有限维向量空间。下一个结果表明这蕴含每个 \(\Delta_i = \mathbb{C}\).

\begin{proposition}\label{pro:18.9}
若 \(\Delta\) 是代数闭域 \(F\) 上的有限维向量空间且是除环，并且 \(F \subseteq Z(\Delta)\)，则 \(\Delta = F\).
\end{proposition}

\begin{proof}
由于 \(F \subseteq Z(\Delta)\)，对每个 \(a \in \Delta\)，由 \(a\) 与 \(F\) 生成的除环是一个域。又由于 \(\Delta\) 在 \(F\) 上有限维，域 \(F(a)\) 是 \(F\) 的有限扩张。因为 \(F\) 代数闭，它没有非平凡有限扩张，故对所有 \(a \in \Delta\) 有 \(F(a) = F\)，即 \(\Delta = F\).
\end{proof}

这个命题证明 \(\mathbb{C}G\) 的 Wedderburn 分解中的每个 \(R_i\) 都是 \(\mathbb{C}\) 上的矩阵环：
\[
\mathbb{C}G \cong M_{n_1}(\mathbb{C}) \times M_{n_2}(\mathbb{C}) \times \cdots \times M_{n_r}(\mathbb{C}).
\]
再由命题 8（5）推出
\[
\sum_{i=1}^{r}n_i^2 = |G|.
\]

本节的最后一个应用是证明 \(r\)（\(\mathbb{C}G\) 的 Wedderburn 分量个数）等于 \(G\) 的共轭类数。为此，先注意命题 8（5）断言 \(r = \dim_{\mathbb C}Z(\mathbb{C}G)\). 我们用另一种方式计算这个维数。

设 \(K_1, \ldots, K_s\) 是 \(G\) 的互不相同的共轭类（回忆它们构成 \(G\) 的一个划分）。对 \(G\) 的每个共轭类 \(K_i\) 令
\[
X_i = \sum_{g \in K_i}g \in \mathbb{C}G.
\]
注意当 \(i \neq j\) 时 \(X_i\) 与 \(X_j\) 没有公共项，故它们是 \(\mathbb{C}G\) 中线性无关的元素。进一步，由于用群元素作共轭置换每个类的元素，有 \(h^{-1}X_ih = X_i\)，即 \(X_i\) 与所有群元素交换。这证明 \(X_i \in Z(\mathbb{C}G)\).

我们证明 \(X_i\) 构成 \(Z(\mathbb{C}G)\) 的一组基，这将证明 \(s = \dim_{\mathbb C}Z(\mathbb{C}G) = r\). 由于 \(X_i\) 线性无关，只需证明它们张成 \(Z(\mathbb{C}G)\). 设 \(X = \sum_{g \in G}a_gg\) 是 \(Z(\mathbb{C}G)\) 的任意元素。由于 \(h^{-1}Xh = X\)，
\[
\sum_{g \in G}a_{hgh^{-1}}gh = \sum_{g \in G}a_gg.
\]
由于 \(G\) 的元素构成 \(\mathbb{C}G\) 的基，上式中 \(g\) 的系数在两边相等：
\[
a_{hgh^{-1}} = a_g.
\]
由于 \(h\) 任意，同一共轭类中所有元素在 \(X\) 中的系数相同，故 \(X\) 可写成 \(X_i\) 的线性组合。

把这些结果总结为下述定理。

\begin{theorem}\label{thm:18.10}
设 \(G\) 是有限群。

（1）\(\mathbb{C}G \cong M_{n_1}(\mathbb{C}) \times M_{n_2}(\mathbb{C}) \times \cdots \times M_{n_r}(\mathbb{C})\).

（2）\(\mathbb{C}G\) 恰有 \(r\) 个互不同构的不可约模，它们的复维数为 \(n_1, n_2, \ldots, n_r\)（从而 \(G\) 恰有 \(r\) 个不等价的不可约复表示，其次数相应为 \(n_1, n_2, \ldots, n_r\)）。

（3）\(\sum_{i=1}^{r}n_i^2 = |G|\).

（4）\(r\) 等于 \(G\) 的共轭类数。
\end{theorem}

\begin{corollary}\label{cor:18.11}
（1）设 \(A\) 是有限阿贝尔群。\(A\) 的每个不可约复表示都是 1 维的（即从 \(A\) 到 \(\mathbb{C}^\times\) 的同态），且 \(A\) 有 \(|A|\) 个不等价的不可约复表示。进一步，\(A\) 的每个有限维复矩阵表示都等价于到对角矩阵群的表示。

（2）任何有限群 \(G\) 的不等价（不可约）1 次复表示个数等于 \(|G:G'|\).
\end{corollary}

\begin{proof}
若 \(A\) 阿贝尔，则 \(\mathbb{C}A\) 是交换环。由于 \(k > 1\) 时 \(k \times k\) 矩阵环不交换，必有每个 \(n_i = 1\). 故 \(r = |A|\)（即 \(A\) 的共轭类数）。由于每个 \(\mathbb{C}A\)-模都是不可约子模的直和，存在一组基，使矩阵关于这组基为对角矩阵。这建立推论的第一部分。

对一般的群 \(G\)，每个 1 次表示 \(\varphi\) 是 \(G\) 到 \(\mathbb{C}^\times\) 的同态，故通过 \(G/G'\) 分解。反之，\(G/G'\) 的每个 1 次表示与自然投射 \(G \to G/G'\) 复合，给出 \(G\) 的 1 次表示。故 \(G\) 的 1 次表示恰是阿贝尔群 \(G/G'\) 的不可约表示。由（1）立刻得到（2）。
\end{proof}

\noindent\textbf{例}

（1）有限阿贝尔群 \(A\) 的不可约复表示（即 \(A\) 到 \(\mathbb{C}^\times\) 的同态）可如下显式描述：把 \(A\) 分解为循环群的直积
\[
A \cong C_1 \times \cdots \times C_n,
\]
其中 \(|C_i| = |\langle x_i\rangle| = d_i\). 把每个 \(x_i\) 映到某个（不必本原的）\(d_i\) 次单位根，并把它延拓到 \(x_i\) 的所有幂以给出同态。由于每个 \(x_i\) 的像有 \(d_i\) 种选择，这个过程定义的 \(A\) 到 \(\mathbb{C}^\times = GL_1(\mathbb{C})\) 的不同同态个数等于 \(|A|\). 由推论 11，这些就是 \(A\) 的全部不可约表示。注意为实现这些表示，域必须包含相应的单位根。下面一个习题探讨循环群在 \(\mathbb{Q}\) 上的不可约表示。

（2）设 \(G = S_3\). 由定理 10，\(G\) 的不可约复表示个数为三（\(S_3\) 的共轭类数）。由于次数的平方和为 6，次数必为 1, 1 与 2. 两个 1 次表示立刻可见：平凡表示，以及由把置换映到其符号（即 \(\sigma\) 为偶置换时 \(\sigma \mapsto +1\)、\(\sigma\) 为奇置换时 \(\sigma \mapsto -1\)）给出的 \(S_3\) 到 \(\{\pm1\}\) 的表示。2 次表示可通过把 3 个基向量上的置换表示（在第 1 节中描述）分解为不可约表示而找到，如下：设 \(S_3\) 通过置换下标作用于向量空间 \(V\) 的基向量 \(e_1, e_2, e_3\) 上。向量 \(t = e_1+e_2+e_3\) 是非零不动向量，故它张成 1 维 \(G\)-不变子空间（这是平凡表示的一个拷贝）。由 Maschke 定理，存在 2 维 \(G\)-不变补空间 \(I\). 注意置换表示不是 1 次表示的和：否则它可由对角矩阵表示，而这些置换在其作用下将彼此交换——这不可能，因为该表示忠实而 \(G\) 非阿贝尔。故 \(I\) 不能进一步分解，于是 \(I\) 提供不可约 2 次表示。事实上，\(I\) 就是第 1 节中描述的「增广」子模：
\[
I = \{w \in V \mid w = a_1e_1+a_2e_2+a_3e_3,\ a_1+a_2+a_3 = 0\}.
\]
显然 \(e_1-e_2\) 与 \(e_2-e_3\) 是 \(I\) 中线性无关的向量，故它们构成这个 2 维空间的一组基。关于 \(I\) 的这组基得到 \(S_3\) 的一个矩阵表示，例如 \(S_3\) 的两个元素上的矩阵表示为
\[
(1\ 2) \mapsto \begin{pmatrix}-1 & 1\\ 0 & 1\end{pmatrix}, \qquad (1\ 2\ 3) \mapsto \begin{pmatrix}0 & -1\\ 1 & -1\end{pmatrix}.
\]

（3）我们分解任意有限群在 \(\mathbb{C}\) 上的正则表示。回忆它是由左 \(\mathbb{C}G\)-模 \(\mathbb{C}G\) 本身提供的表示。由定理 10，\(\mathbb{C}G\) 首先是双边理想的直积：
\[
\mathbb{C}G \cong M_{n_1}(\mathbb{C}) \times M_{n_2}(\mathbb{C}) \times \cdots \times M_{n_r}(\mathbb{C}).
\]
现在由命题 6（4），每个 \(M_{n_i}(\mathbb{C})\) 进一步分解为 \(n_i\) 个同构单左理想的直和。这些左理想给出不可约 \(\mathbb{C}G\)-模的一组完全同构类代表。故 \(G\) 的正则表示（在 \(\mathbb{C}\) 上）分解为 \(G\) 的所有不可约表示之直和，其中每个出现的重数等于该不可约表示的次数。

我们记录 \(\mathbb{C}G\) 的一个附加性质，其证明见第 19.2 节。

\begin{theorem}\label{thm:18.12}
有限群 \(G\) 的每个不可约复表示的次数整除 \(G\) 的阶，即按定理 10 的记号，对 \(i = 1, 2, \ldots, r\) 有 \(n_i\) 整除 \(|G|\).
\end{theorem}

下一节我们将用群元素描述 \(\mathbb{C}G\) 的本原中心幂等元。

\subsection*{习题}

设 \(G\) 是有限群，\(R\) 是含 1 的环。

\begin{enumerate}
\item 证明 Wedderburn 定理的条件（1）与（2）等价。

\item 证明 Wedderburn 定理中（3）蕴含（2）。〔设 \(Q\) 是 \(R\)-模 \(N\) 的子模。用 Zorn 引理证明存在关于 \(Q \cap M = 0\) 极大的子模 \(M\). 若 \(Q+M = N\)，则（2）成立；否则设 \(M_1\) 是 \(N/M\) 中某个不含于 \((Q+M)/M\) 的单模在 \(N\) 中的完全原像，并论证 \(M_1\) 与 \(M\) 的极大性矛盾。〕

\item 证明 Wedderburn 定理中（4）蕴含（3）。〔设 \(N\) 是非零 \(R\)-模。先通过考虑循环子模证明 \(N\) 含单子模。再对单子模的直和之集（适当排序后）用 Zorn 引理，证明 \(N\) 含极大完全可约子模 \(M\). 若 \(M \neq N\)，设 \(M_1\) 是 \(N/M\) 中某个单模在 \(N\) 中的完全原像，并与 \(M\) 的极大性矛盾。〕

\item 证明 Wedderburn 定理中（5）蕴含（4）。〔用命题 6 与命题 8 的证明方法把每个 \(R_i\) 分解为左 \(R\)-模。〕
\end{enumerate}

下面六个习题建立关于环与模的一些一般结果，它们蕴含 Wedderburn 定理剩下的一个方向：由（2）推出（5）。这些习题中假设 \(R\) 满足（2）：每个 \(R\)-模都是内射的。

\begin{enumerate}
\setcounter{enumi}{4}
\item 证明 \(R\) 在左理想上满足降链条件（D.C.C.）。推出 \(R\) 是有限多个左理想的直和。〔若不是，则证明作为左 \(R\)-模，\(R\) 是无限多个非零子模的直和。把这个直和中的元素 1 写出来以导出矛盾。〕

\item 证明 \(R = R_1 \times R_2 \times \cdots \times R_r\)，其中 \(R_j\) 是双边理想且是单环（即没有真非零双边理想）。证明每个 \(R_j\) 有单位元，且在左理想上满足 D.C.C.〔用上一习题证明 \(R\) 有极小双边理想 \(R_1\). 作为左 \(R\)-模有 \(R = R_1 \oplus R'\)，其中 \(R'\) 是某个左理想。证明 \(R'\) 是右理想，并用 D.C.C. 归纳进行。〕

\item 设 \(S\) 是含 1 的单环，在左理想上满足 D.C.C.，\(L\) 是 \(S\) 的极小左理想。证明作为左 \(S\)-模有 \(S \cong L \oplus \cdots \oplus L\)（\(n\) 项）。〔由单性论证 \(LS = S\)，故对最小的 \(n\) 有 \(1 = l_1s_1 + \cdots + l_ns_n\)，其中 \(l_i \in L\)、\(s_i \in S\). 证明映射 \((x_1, \ldots, x_n) \mapsto x_1s_1 + \cdots + x_ns_n\) 是左 \(S\)-模的满同态；用 \(L\) 与 \(n\) 的极小性证明它是单射。〕

\item 设 \(A\) 是任意含 1 的环，\(L\) 是任意左 \(A\)-模，\(P\) 是 \(n\) 个 \(L\) 的直和。

（a）证明环同构 \(\operatorname{Hom}_A(P, P) \cong M_n(D)\)，其中 \(D = \operatorname{Hom}_A(L, L)\)（环 \(\operatorname{Hom}_A(X, X)\) 中的乘法是复合，参看第 10.2 节命题 2（4））。

（b）推出若 \(L\) 是单 \(A\)-模，则 \(\operatorname{Hom}_A(P, P)\) 同构于除环上的矩阵环。〔用 Schur 引理与（a）。〕

（c）证明环同构 \(\operatorname{Hom}_A(A, A) \cong A^{opp}\)，其中 \(A^{opp}\) 是 \(A\) 的反环（元素与加法与 \(A\) 中相同，但乘积 \(x \cdot y\) 在 \(A^{opp}\) 中的值是 \(A\) 中计算的 \(yx\)），参看第 17.4 节末。〔任何同态由它在 1 上的值决定。〕

\item 证明若 \(S\) 是含 1 的单环且在左理想上满足 D.C.C.，则 \(S \cong M_n(\Delta)\)，其中 \(\Delta\) 是某个除环。（这个结果连同习题 6 完成了「（2）蕴含（5）」证明的存在性部分。）〔用习题 7 与 8 证明 \(S^{opp} \cong \operatorname{Hom}_S(P, P) \cong M_n(D)\)，其中 \(D\) 是某个除环。然后证明 \(S \cong M_n(\Delta)\)，其中 \(\Delta\) 是除环 \(D^{opp}\). 这里 \(P\) 是习题 7 中的左 \(S\)-模。〕

\item 证明上一习题的同构 \(S \cong M_n(\Delta)\) 中的 \(\Delta\) 与 \(n\) 由 \(S\) 唯一确定（这证明 Wedderburn 定理中的唯一性陈述），如下进行。假设作为环有 \(S = M_n(\Delta) \cong M_{n'}(\Delta')\)，其中 \(\Delta\) 与 \(\Delta'\) 是除环。

（a）证明 \(\Delta \cong \operatorname{Hom}_S(L, L)\)，其中 \(L\) 是 \(S\) 的极小左理想。推出 \(\Delta \cong \Delta'\).〔用命题 6（4）。〕

（b）证明除环 \(\Delta\) 上的有限生成（左）模有「基」（线性无关的生成集），且任意两组基有相同基数。推出 \(n = n'\).〔模仿第 11.1 节推论 4（2）的证明。〕

\item 证明若 \(R\) 是含 1 的环，且每个 \(R\)-模都是自由模，则 \(R\) 是除环。

\item 设 \(F\) 是域，\(f(x) \in F[x]\)，\(R = F[x]/(f(x))\). 求 \(f(x)\) 在 \(F[x]\) 中的分解应满足的充要条件，使 \(R\) 是半单环。当 \(R\) 半单时，描述它的 Wedderburn 分解。〔参见第 9.5 节命题 16。〕

\item 设 \(G\) 是 \(n\) 阶循环群，\(R = \mathbb{Q}G\). 描述 \(R\) 的 Wedderburn 分解，并求出 \(G\) 在 \(\mathbb{Q}\) 上不可约表示的个数与次数。特别地，证明若 \(n = p\) 是素数，则 \(G\) 在 \(\mathbb{Q}\) 上恰有一个非平凡不可约表示，且其次数为 \(p-1\).〔回忆第 1 节第一个例子中 \(\mathbb{Q}G \cong \mathbb{Q}[x]/(x^n-1)\). 用第 9.5 节命题 16 与第 13.6 节的结果。〕

\item 设 \(p\) 是素数，\(F\) 是 \(p\) 元域，\(G\) 是 3 阶循环群，\(R = FG\). 对 \(p = 2\) 与 \(p = 7\) 的每一种情形，描述 \(R\) 的 Wedderburn 分解，并求出 \(G\) 在 \(F\) 上不可约表示的个数与次数。

\item 证明若 \(P\) 是某个素数 \(p\) 的 \(p\)-群，则 \(P\) 有忠实的不可约复表示当且仅当 \(Z(P)\) 循环。〔用第 1 节习题 18、定理 6.1（2）与例 3。〕

\item 证明若 \(V\) 是不可约 \(FG\)-模且 \(F\) 是代数闭域，则 \(\operatorname{Hom}_{FG}(V, V)\) 同构于 \(F\)（作为环）。

\item 设 \(F\) 是域，\(R = M_n(F)\)，\(M\) 是唯一的不可约 \(R\)-模。证明 \(\operatorname{Hom}_R(M, M)\) 同构于 \(F\)（作为环）。

\item 求出 \(M_n(\Delta)\) 的所有双边理想。
\end{enumerate}

\section{特征标理论与正交关系}

一般地，对于阶很大的群，表示难以计算，也（即使并非不可能写出）难以驾驭。例如，一个 100 次的矩阵表示涉及含 10000 个元素的矩阵，而为了描述这个表示（即使在群的一组生成元上）可能需要许多个 \(100 \times 100\) 矩阵。然而，有一些引人注目的例子，其中高次表示已被计算出来并得到有效使用。一个实例是 Z. Janko 在 1965 年对单群 \(J_1\) 的构造（单群的存在性问题已在第 6.2 节末讨论）。Janko 在研究单群的某些性质时发现，若某个单群具有这些性质，则它必有阶 175,560，且由两个元素生成。此外，他证明具有这些性质的假想单群必有一个 \(\mathbb{F}_{11}\) 上的 7 维表示，其两个生成元映到两个矩阵
\[
\begin{pmatrix}
0&1&0&0&0&0&0\\
0&0&1&0&0&0&0\\
0&0&0&1&0&0&0\\
0&0&0&0&1&0&0\\
0&0&0&0&0&0&-3\\
0&0&0&0&0&0&1\\
1&0&0&0&0&0&0
\end{pmatrix}
\qquad\text{与}\qquad
\begin{pmatrix}
-3&2&-1&-1&-3&-1&-3\\
-2&1&1&3&1&3&3\\
-1&-1&-3&-1&-3&-3&2\\
-1&-3&-1&-3&-3&2&-1\\
-1&-3&-3&2&-1&-1&1\\
1&3&3&-2&1&1&3\\
3&3&-2&1&1&3&1
\end{pmatrix}
\]
（扫描件中这两个矩阵的元素辨识度较低，以上按所见字形转录，建议对照原书核对。）

（注意对任何单群 \(S\)，\(S\) 到 \(GL_n(F)\) 的、不把所有群元素映到单位矩阵的表示都是忠实表示，故 \(S\) 同构于它在 \(GL_n(F)\) 中的像。）特别地，Janko 的计算表明满足他的性质的单群若存在则唯一。M. Ward 证明了这两个矩阵确实生成 \(GL_7(\mathbb{F}_{11})\) 的一个 175,560 阶子群，从而确实存在满足 Janko 性质的单群。

类似地，S. Norton、R. Parker 与 J. Thackray 用 \(\mathbb{F}_2\) 上的一个 112 维表示构造了 86,775,571,046,077,562,880 阶单群 \(J_4\). 这个群被证明由两个元素生成，而在他们的分析过程中计算出了这两个生成元在 \(GL_{112}(\mathbb{F}_2)\) 中的显式矩阵。

1981 年，R. Griess 构造了最大的零星群，即所谓 \textbf{Monster}，其阶为
\[
2^{46}\cdot 3^{20}\cdot 5^9\cdot 7^6\cdot 11^2\cdot 13^3\cdot 17\cdot 19\cdot 23\cdot 29\cdot 31\cdot 41\cdot 47\cdot 59\cdot 71.
\]
他的证明涉及计算 \(\mathbb{C}\) 上一个 196,884 维代数的自同构，并由此给出 Monster 的一个该次数的表示。

类似地，一般难以写出大次数 \(n\) 的置换表示 \(\varphi : G \to S_n\) 所对应的显式置换。然而，有一些数值不变量，例如置换的符号与轮换型 \(\pi(g)\)，而这些数值不变量可能比置换本身更容易计算（即可能不必真正写出置换本身就能确定元素的轮换型，如第 14.8 节中 \(\mathbb{Q}\) 上伽罗瓦群的计算）。仅这些不变量在给定情形下就可能提供足够的信息来完成某些分析，例如证明一个给定的群不单（如第 6.2 节所示）。此外，刚刚提到的这些不变量不依赖于集合 \(\{1, 2, \ldots, n\}\) 的标号（即它们不依赖于 \(S_n\) 中的「基变换」），并且对于在 \(G\) 中共轭的元素它们相同。

本节我们说明如何把数值不变量联系到线性表示上。这些不变量只依赖于表示的等价类（同构型）。换句话说，对每个表示 \(\varphi : G \to GL_n(F)\)，我们将把 \(F\) 的一个元素联系到每个矩阵 \(\varphi(g)\) 上，并将看到这个数在许多情形下可以不必知道矩阵 \(\varphi(g)\) 就算出。此外，我们将看到这些不变量不依赖于 \(\varphi\) 的相似类（即若把表示 \(\varphi\) 换成等价表示，则对固定的 \(g \in G\) 这些数不变），并且在某种意义上刻画了 \(G\) 的表示的相似类。

本节中 \(G\) 始终是有限群，而眼下 \(F\) 是任意域。所考虑的所有表示都假定是有限次的。

\begin{definition}
（1）\textbf{类函数}是从 \(G\) 到 \(F\) 的、在 \(G\) 的共轭类上取常值的任何函数，即 \(f : G \to F\)，使对所有 \(g, x \in G\) 有 \(f(g^{-1}xg) = f(x)\).

（2）若 \(\varphi\) 是由 \(FG\)-模 \(V\) 提供的 \(G\) 的表示，则 \(\varphi\) 的\textbf{特征标}是由 \(\chi(g) = \operatorname{tr}\varphi(g)\) 定义的函数 \(\chi : G \to F\)，其中 \(\operatorname{tr}\varphi(g)\) 是 \(\varphi(g)\) 关于 \(V\) 的某组基的矩阵的迹（即该矩阵对角元素之和）。按表示是否不可约或可约，称特征标是\textbf{不可约}或\textbf{可约}的。特征标的\textbf{次数}是提供它的任何表示的次数。
\end{definition}

在这个定义的第二部分的记号下，我们也把 \(\chi\) 称为由 \(FG\)-模 \(V\) 提供的特征标。一般地，特征标不是从群到域的加法群或乘法群的同态。

\noindent\textbf{例}

（1）平凡表示的特征标是函数 \(\chi(g) = 1\)（对所有 \(g \in G\)）。这个特征标称为 \(G\) 的\textbf{主特征标}。

（2）对于 1 次表示，特征标与表示通常等同（把 \(1 \times 1\) 矩阵与它的元素等同）。故对阿贝尔群，不可约复表示与它们的特征标是同一回事（参看推论 11）。

（3）设 \(\pi : G \to S_n\) 是置换表示，\(\varphi\) 是向量空间 \(V\) 的基 \(e_1, \ldots, e_n\) 上由此得到的线性表示：
\[
\varphi(g)(e_i) = e_{\pi(g)(i)}
\]
（参看第 1 节例 4）。关于这组基，若 \(\pi(g)\) 固定 \(i\)，则 \(\varphi(g)\) 的矩阵在对角位置 \(i,i\) 处为 1；否则 \(\varphi(g)\) 的矩阵在位置 \(i,i\) 处为零。故若 \(\eta\) 是 \(\varphi\) 的特征标，则
\[
\eta(g) = g \text{ 在 } \{1, 2, \ldots, n\} \text{ 上的不动点个数}.
\]
特别地，若 \(\pi\) 是由左乘作用在 \(G\) 的某个子群 \(H\) 的左陪集之集上所得的置换表示，则所得特征标称为 \(G\) 在 \(H\) 上的\textbf{置换特征标}。

（4）例 3 中 \(\pi\) 为 \(G\) 的正则置换表示的特殊情形值得记录：若 \(\varphi\) 是 \(G\) 的正则表示（由模 \(FG\) 提供）而 \(\rho\) 是它的特征标，则
\[
\rho(g) = \begin{cases}0, & \text{若 } g \neq 1\\ |G|, & \text{若 } g = 1.\end{cases}
\]
\(G\) 的正则表示的特征标称为 \(G\) 的\textbf{正则特征标}。注意这提供了特征标取值为 0、且不是从 \(G\) 到 \(F\) 或 \(F^\times\) 的群同态的具体例子。

（5）设 \(\varphi : D_{2n} \to GL_2(\mathbb{R})\) 是第 1 节第二组例子中例 6 描述的显式矩阵表示。若 \(\chi\) 是 \(\varphi\) 的特征标，则由给定 \(2 \times 2\) 矩阵的迹可知 \(\chi(r) = 2\cos(2\pi/n)\) 且 \(\chi(s) = 0\). 由于 \(\varphi\) 把 \(D_{2n}\) 的单位元映到 \(2 \times 2\) 单位矩阵，\(\chi(1) = 2\).

（6）设 \(\varphi : Q_8 \to GL_2(\mathbb{C})\) 是第 1 节第二组例子中例 7 描述的显式矩阵表示。若 \(\chi\) 是 \(\varphi\) 的特征标，则由给定 \(2 \times 2\) 矩阵的迹可知 \(\chi(i) = 0\) 且 \(\chi(j) = 0\). 由于元素 \(-1 \in Q_8\) 映到负的 \(2 \times 2\) 单位矩阵，\(\chi(-1) = -2\). 由于 \(\varphi\) 把 \(Q_8\) 的单位元映到 \(2 \times 2\) 单位矩阵，\(\chi(1) = 2\).

（7）设 \(\varphi : Q_8 \to GL_4(\mathbb{R})\) 是第 1 节第二组例子中例 8 描述的矩阵表示。若 \(\chi\) 是 \(\varphi\) 的特征标，则通过观察所给出的矩阵可知 \(\chi(i) = \chi(j) = 0\). 由于 \(\varphi\) 把 \(Q_8\) 的单位元映到 \(4 \times 4\) 单位矩阵，\(\chi(1) = 4\).

对 \(n \times n\) 矩阵 \(A\) 与 \(B\)，直接计算表明 \(\operatorname{tr}AB = \operatorname{tr}BA\). 若 \(A\) 可逆，则推出
\[
\operatorname{tr}A^{-1}BA = \operatorname{tr}B.
\]
故表示的特征标不依赖于提供它的向量空间的基的选取，即
\[
\text{等价表示有相同的特征标}. \tag{18.1}
\]

设 \(\varphi\) 是 \(G\) 的次数为 \(n\) 的表示，特征标为 \(\chi\). 由于对所有 \(g, x \in G\) 有 \(\varphi(g^{-1}xg) = \varphi(g)^{-1}\varphi(x)\varphi(g)\)，取迹表明
\[
\text{表示的特征标是类函数}. \tag{18.2}
\]
由于 \(n \times n\) 单位矩阵的迹为 \(n\)，而 \(\varphi\) 把 \(G\) 的单位元映到恒等线性变换（或矩阵），
\[
\chi(1) \text{ 是 } \varphi \text{ 的次数}. \tag{18.3}
\]

若 \(V\) 是所提供的表示具有特征标 \(\chi\) 的 \(FG\)-模，则群环 \(FG\) 的每个元素都作为从 \(V\) 到 \(V\) 的线性变换作用。故每个 \(\sum_{g \in G}a_gg \in FG\) 在看作从 \(V\) 到 \(V\) 的线性变换时有迹。\(g \in G\) 作用在 \(V\) 上的迹按定义是 \(\chi(g)\). 由于任何矩阵线性组合的迹是迹的线性组合，\(\sum_{g \in G}a_gg\) 作用在 \(V\) 上的迹是 \(\sum_{g \in G}a_g\chi(g)\). 注意这个 \(FG\) 上的迹函数是 \(G\) 的特征标 \(\chi\) 到从 \(FG\) 到 \(F\) 的 \(F\)-线性变换的唯一延拓。这样我们将把 \(G\) 的特征标也看作定义在群环 \(FG\) 上。

注意上面例 3 中，若域 \(F\) 的特征 \(p > 0\)，则特征标模 \(p\) 的值可能为零，即使不动点个数不为零。为了绕开这类异常，并利用在 \(F\) 代数闭时所得的 Wedderburn 定理的推论，我们再次把域特殊化到复数（或任何特征为 0 的代数闭域）。由前一节的结果
\[
\mathbb{C}G \cong M_{n_1}(\mathbb{C}) \times M_{n_2}(\mathbb{C}) \times \cdots \times M_{n_r}(\mathbb{C}). \tag{18.4}
\]

本节余下部分固定下述记号：
\[
M_1, M_2, \ldots, M_r \text{ 是不等价的不可约 } \mathbb{C}G\text{-模}, \tag{18.5}
\]
\[
\chi_i \text{ 是由 } M_i \text{ 提供的特征标}, \quad 1 \leq i \leq r.
\]
故 \(r\) 是 \(G\) 的共轭类数；必要时可重新标号 \(M_1, \ldots, M_r\)，使 \(\chi_i\) 的次数对所有 \(i\) 都是 \(n_i\)（这也就是 \(M_i\) 在 \(\mathbb{C}\) 上的维数）。

现在每个（有限维）\(\mathbb{C}G\)-模 \(M\) 都同构（等价）于不可约模的直和：
\[
M \cong a_1M_1 \oplus a_2M_2 \oplus \cdots \oplus a_rM_r, \tag{18.6}
\]
其中 \(a_i\) 是非负整数，表示不可约模 \(M_i\) 在这个直和分解中的重数，即
\[
a_iM_i = \underbrace{M_i \oplus M_i \oplus \cdots \oplus M_i}_{a_i \text{ 次}}.
\]

注意若表示由模 \(M\) 提供，而 \(M = M_1 \oplus M_2\)，则我们可以取 \(M\) 的一组基，它由 \(M_1\) 的一组基与 \(M_2\) 的一组基合并而成。关于这组基的矩阵表示为
\[
\varphi(g) = \begin{pmatrix}\varphi_1(g) & 0\\ 0 & \varphi_2(g)\end{pmatrix},
\]
其中 \(\varphi_i\) 是由 \(M_i\) 提供的表示，\(i = 1, 2\). 立刻可以看出，若 \(\psi\) 是 \(\varphi\) 的特征标而 \(\psi_i\) 是 \(\varphi_i\) 的特征标，则 \(\psi(g) = \psi_1(g)+\psi_2(g)\)，即 \(\psi = \psi_1+\psi_2\). 由归纳得到：
\[
\text{表示的特征标是直和分解中出现的各成分的特征标之和}. \tag{18.7}
\]

若 \(\psi\) 是上面（6）中模 \(M\) 提供的特征标，则这给出
\[
\psi = a_1\chi_1+a_2\chi_2+\cdots+a_r\chi_r. \tag{18.8}
\]
故每个（复）特征标都是不可约（复）特征标的非负整数线性组合。反之，通过取模的直和可看出，每个这样的特征标之和都是 \(G\) 的某个复表示的特征标。

接下来我们证明特征标与复表示的等价类之间的对应是双射。设 \(z_1, z_2, \ldots, z_r\) 是前一节描述的 \(\mathbb{C}G\) 的本原中心幂等元。由于它们正交（或等价地，由于它们是 \(\mathbb{C}G\) 分解为 \(r\) 个「在一处为 1、其余处为零」的含 1 子环之直积中的 \(r\) 元组），\(z_1, \ldots, z_r\) 是 \(\mathbb{C}G\) 中 \(\mathbb{C}\)-线性无关的元素。如上，每个不可约特征标 \(\chi_i\) 是 \(\mathbb{C}G\) 上的函数。由命题 8（3）我们有

（a）若 \(j \neq i\) 则 \(z_jM_i = 0\)，即 \(z_j\) 在 \(M_i\) 上作为零矩阵作用，故 \(\chi_i(z_j) = 0\)；以及

（b）\(z_i\) 在 \(M_i\) 上作为恒等作用，故 \(\chi_i(z_i) = n_i\).

故 \(\chi_1, \ldots, \chi_r\) 是独立集合 \(z_1, \ldots, z_r\) 的对偶基的倍数，从而是线性无关的函数。现在，若上面（6）中描述的 \(\mathbb{C}G\)-模 \(M\) 可以按另一种方式分解为不可约模，比如
\[
M \cong b_1M_1 \oplus b_2M_2 \oplus \cdots \oplus b_rM_r,
\]
则我们得到关系式
\[
a_1\chi_1+a_2\chi_2+\cdots+a_r\chi_r = b_1\chi_1+b_2\chi_2+\cdots+b_r\chi_r.
\]
由不可约特征标的线性无关性，对所有 \(i \in \{1, \ldots, r\}\) 有 \(b_i = a_i\). 故在 \(M\) 的任何不可约直和分解中，不可约模 \(M_i\) 的重数相同，\(1 \leq i \leq r\). 特别地，
\[
\text{两个表示等价当且仅当它们有相同的特征标}. \tag{18.9}
\]

这个唯一性也可以用另一种方式看出。首先用命题 8（2）把任意有限维 \(\mathbb{C}G\)-模 \(M\) 唯一分解为
\[
M = z_1M \oplus z_2M \oplus \cdots \oplus z_rM.
\]
由同一命题的（4），\(z_iM\) 是单模的直和，其中每一个都同构于 \(M_i\). 通过计算维数可知，\(M_i\) 在 \(M\) 的直和分解中的重数等于 \(\dim z_iM / \dim_{\mathbb C}M_i\). 这证明 \(M_i\) 在 \(M\) 的任何单子模直和分解中的重数被唯一确定。

注意，正如 \(F[x]\)-模分解为循环子模时一样，一个 \(\mathbb{C}G\)-模可能有许多不可约直和分解——只有重数是唯一的（另见习题）。更确切地说，与单个线性变换的 Jordan 标准形作比较，直和项 \(a_iM_i = M_i \oplus \cdots \oplus M_i\)（\(a_i\) 次）等于子模 \(z_iM\)，它是单个特征值所对应广义特征空间的类比。\(M\) 的这个子模是唯一的（正如广义特征空间一样），称为 \(M\) 的 \(i\) 型\textbf{等型成分}。在等型成分内部，直和项 \(M_i\) 类似于 1 维特征空间；正如自同态的特征空间没有唯一的一组基，等型成分也没有唯一的一组基。若 \(G = \langle g\rangle\) 是有限循环群，则 \(G\) 的等型成分与 \(g\) 的广义特征空间相同。

注意全体（复值）类函数构成的向量空间有一组基，由在给定类上取 1、在其他所有类上取 0 的函数组成。这样的函数有 \(r\) 个，其中 \(r\) 是 \(G\) 的共轭类数，故类函数的复向量空间的维数为 \(r\). 由于 \(G\) 的（复）不可约特征标个数等于共轭类数，且这些特征标是线性无关的类函数，我们看到
\[
\text{不可约特征标是全体复类函数空间的一组基}. \tag{18.10}
\]

特征标理论的下一步是在类函数空间上赋予 Hermite 内积结构，并证明不可约特征标关于这个内积构成一组标准正交基。对类函数 \(\theta\) 与 \(\psi\) 定义
\[
(\theta, \psi) = \frac{1}{|G|}\sum_{g \in G}\theta(g)\overline{\psi(g)}
\]
（其中横线表示复共轭）。容易验证 \((\ ,\ )\) 是 Hermite 的，即对 \(\alpha, \beta \in \mathbb{C}\)：

（a）\((\alpha\theta_1+\beta\theta_2, \psi) = \alpha(\theta_1, \psi)+\beta(\theta_2, \psi)\)，

（b）\((\theta, \alpha\psi_1+\beta\psi_2) = \bar\alpha(\theta, \psi_1)+\bar\beta(\theta, \psi_2)\)，以及

（c）\((\theta, \psi) = \overline{(\psi, \theta)}\).

我们的主要目标是证明不可约特征标关于这个 Hermite 形式构成复类函数空间的一组标准正交基（我们已知它们是一组基）。一旦在下一个命题中显式确定了本原中心幂等元，这个事实就将由这些幂等元的正交性推出。

\begin{proposition}\label{pro:18.13}
设 \(z_1, \ldots, z_r\) 是 \(\mathbb{C}G\) 中正交本原中心幂等元，其标号使得 \(z_i\) 在不可约 \(\mathbb{C}G\)-模 \(M_i\) 上作为恒等作用，而 \(\chi_i\) 是由 \(M_i\) 提供的特征标。则
\[
z_i = \frac{\chi_i(1)}{|G|}\sum_{g \in G}\chi_i(g^{-1})g.
\]
\end{proposition}

\begin{proof}
设 \(z = z_i\)，写成
\[
z = \sum_{g \in G}a_gg.
\]
回忆本节例 4 中 \(\rho\) 是 \(G\) 的正则特征标，则
\[
\rho(g) = \begin{cases}0, & \text{若 } g \neq 1\\ |G|, & \text{若 } g = 1,\end{cases} \tag{18.11}
\]
并回忆本节例 2 中
\[
\rho = \sum_{j=1}^{r}\chi_j(1)\chi_j. \tag{18.12}
\]
为求系数 \(a_g\)，把 \(\rho\) 作用在 \(zg^{-1}\) 上，并用 \(\rho\) 的线性性以及等式（11）得到
\[
\rho(zg^{-1}) = a_g|G|.
\]
再用（12）计算 \(\rho(zg^{-1})\) 给出
\[
\sum_{j=1}^{r}\chi_j(1)\chi_j(zg^{-1}) = a_g|G|. \tag{18.13}
\]
设 \(\varphi_j\) 是由 \(M_j\) 提供的表示，\(1 \leq j \leq r\). 由于 \(zg^{-1}\) 在 \(M_j\) 上作用为 \(\varphi_j(z)\varphi_j(g^{-1})\)，而我们已经观察到当 \(j \neq i\) 时 \(\varphi_j(z) = 0\)、当 \(j = i\) 时 \(\varphi_i(z)\) 是 \(M_i\) 上的恒等自同态，故
\[
\chi_j(zg^{-1}) = \begin{cases}0, & j \neq i\\ \chi_i(g^{-1}), & j = i.\end{cases}
\]
这证明 \(\chi_j(zg^{-1}) = \chi_i(g^{-1})\delta_{ij}\)，其中 \(\delta_{ij}\) 当 \(i \neq j\) 时为零、当 \(i = j\) 时为 1（称为 \textbf{Kronecker delta}）。把它代入等式（13）给出
\[
a_g = \frac{1}{|G|}\chi_i(1)\chi_i(g^{-1}).
\]
这就是命题所述式子中 \(g\) 的系数，从而完成了证明。
\end{proof}

不可约特征标的标准正交性将直接由中心本原幂等元的正交性经下述计算推出：
\[
z_iz_j = \delta_{ij}z_i = \frac{\chi_i(1)\chi_j(1)}{|G|^2}\sum_{g,h \in G}\chi_i(g^{-1})\chi_j(h^{-1})gh
= \frac{\chi_i(1)\chi_j(1)}{|G|^2}\sum_{y \in G}\left(\sum_{x \in G}\chi_i(xy^{-1})\chi_j(x^{-1})\right)y,
\]
（为从前者得到后一个和，用 \(y\) 代替 \(gh\)，用 \(x\) 代替 \(h\)）。由于 \(G\) 的元素构成 \(\mathbb{C}G\) 的基，我们可以把系数与命题 13 中求得的 \(z_i\) 的系数比较，得到（\(g\) 的系数满足）
\[
\delta_{ij}\frac{\chi_i(1)}{|G|}\chi_i(g^{-1}) = \frac{\chi_i(1)\chi_j(1)}{|G|^2}\sum_{x \in G}\chi_i(xg^{-1})\chi_j(x^{-1}).
\]
化简（并把 \(g\) 换成 \(g^{-1}\)）给出
\[
\delta_{ij}\chi_i(g) = \frac{\chi_j(1)}{|G|}\sum_{x \in G}\chi_i(xg)\chi_j(x^{-1}) \qquad \text{对所有 } g \in G. \tag{18.14}
\]
在（14）中取 \(g = 1\) 给出
\[
\delta_{ij}\chi_i(1) = \frac{\chi_j(1)}{|G|}\sum_{x \in G}\chi_i(x)\chi_j(x^{-1}). \tag{18.15}
\]
右端的和若把 \(\chi_j(x^{-1})\) 换成它的复共轭 \(\overline{\chi_j(x)}\)，就恰好是内积 \((\chi_i, \chi_j)\) 的 \(|G|\) 倍；而这一步正是下一个命题的内容。

\begin{proposition}\label{pro:18.14}
若 \(\psi\) 是 \(G\) 的任意特征标，则 \(\psi(x)\) 是 \(\mathbb{C}\) 中单位根之和，且对所有 \(x \in G\) 有 \(\psi(x^{-1}) = \overline{\psi(x)}\).
\end{proposition}

\begin{proof}
设 \(\varphi\) 是特征标为 \(\psi\) 的表示，固定元素 \(x \in G\)，并设 \(|x| = k\). 由于 \(\varphi(x)\) 的极小多项式整除 \(X^k-1\)（故有互不相同的根），存在底空间的一组基，使 \(\varphi(x)\) 关于这组基的矩阵是对角矩阵，且对角线上是 \(k\) 次单位根。由于 \(\psi(x)\) 是对角元素之和（且不依赖于基的选取），\(\psi(x)\) 是单位根之和。进一步，若 \(\varepsilon\) 是单位根，则 \(\varepsilon^{-1} = \bar\varepsilon\). 故对角线上为单位根的对角矩阵之逆是对角线上为这些单位根的复共轭的对角矩阵。由于和的复共轭是复共轭之和，
\[
\psi(x^{-1}) = \operatorname{tr}\varphi(x^{-1}) = \sum_i\varepsilon_i^{-1} = \sum_i\bar\varepsilon_i = \overline{\operatorname{tr}\varphi(x)} = \overline{\psi(x)}.
\]
\end{proof}

请记住，在命题 14 的证明中我们先固定一个群元素 \(x\)，然后再选取表示空间的一组基使 \(\varphi(x)\) 为对角矩阵。始终可以把单个元素对角化，但可以同时把所有 \(\varphi(x)\) 对角化当且仅当 \(\varphi\) 相似于 1 次表示之和。

把上述命题与等式（15）结合即证明：

\begin{theorem}\label{thm:18.15}
（\textbf{群特征标的第一正交关系}）设 \(G\) 是有限群，\(\chi_1, \ldots, \chi_r\) 是 \(G\) 在 \(\mathbb{C}\) 上的不可约特征标。则关于上述内积 \((\ ,\ )\) 有
\[
(\chi_i, \chi_j) = \delta_{ij},
\]
且不可约特征标是类函数空间的一组标准正交基。特别地，若 \(\theta\) 是任意类函数，则
\[
\theta = \sum_{i=1}^{r}(\theta, \chi_i)\chi_i.
\]
\end{theorem}

\begin{proof}
我们刚刚建立了不可约特征标是类函数空间的一组标准正交基。若 \(\theta\) 是任意类函数，写成 \(\theta = \sum_{i=1}^{r}a_i\chi_i\)（\(a_i \in \mathbb{C}\)），则由 Hermite 积的线性性有 \(a_i = (\theta, \chi_i)\)，即所述结论。
\end{proof}

我们不加证明地列出第二正交关系；本书中的应用不需要它。

\begin{theorem}\label{thm:18.16}
（\textbf{群特征标的第二正交关系}）在上述记号下，对任意 \(x, y \in G\)，
\[
\sum_{i=1}^{r}\chi_i(x)\overline{\chi_i(y)} = \begin{cases}|C_G(x)|, & \text{若 } x \text{ 与 } y \text{ 在 } G \text{ 中共轭},\\ 0, & \text{否则}.\end{cases}
\]
\end{theorem}

\begin{definition}
对 \(G\) 上任意类函数 \(\theta\)，\(\theta\) 的\textbf{范数}定义为 \(\|\theta\| = (\theta, \theta)^{1/2}\).
\end{definition}

当类函数按不可约特征标写出，\(\theta = \sum a_i\chi_i\) 时，其范数容易算得为
\[
\|\theta\|^2 = (\theta, \theta) = \sum_{i=1}^{r}|a_i|^2.
\]
由此推出：

一个特征标的范数为 1 当且仅当它不可约。

最后注意，特征标 \(\theta\) 与 \(\psi\) 的内积的计算可如下简化。若 \(K_1, \ldots, K_r\) 是 \(G\) 的共轭类，大小分别为 \(d_1, \ldots, d_r\)，代表元分别为 \(g_1, \ldots, g_r\)，则值 \(\theta(g_i)\overline{\psi(g_i)}\) 在 \((\theta, \psi)\) 的和中出现 \(d_i\) 次（对应 \(K_i\) 的每个元素一次）。收集这些项给出
\[
(\theta, \psi) = \frac{1}{|G|}\sum_{i=1}^{r}d_i\theta(g_i)\overline{\psi(g_i)},
\]
一个只对 \(G\) 的共轭类代表元求和的式子。特别地，\(\theta\) 的范数由
\[
\|\theta\|^2 = (\theta, \theta) = \frac{1}{|G|}\sum_{i=1}^{r}d_i|\theta(g_i)|^2
\]
给出。

\noindent\textbf{例}

（1）设 \(G = S_3\)，\(\eta\) 是本节开头例子中描述的 3 次置换特征标。回忆 \(\eta(\sigma)\) 等于 \(\{1,2,3\}\) 中被 \(\sigma\) 固定的元素个数。\(S_3\) 的共轭类由 \(1\)、\((1\ 2)\) 与 \((1\ 2\ 3)\) 代表，大小分别为 1、3、2，且 \(\eta(1) = 3\)、\(\eta((1\ 2)) = 1\)、\(\eta((1\ 2\ 3)) = 0\). 故
\[
\|\eta\|^2 = \frac{1}{6}\left[1\cdot\eta(1)^2+3\cdot\eta((1\ 2))^2+2\cdot\eta((1\ 2\ 3))^2\right] = \frac{1}{6}(9+3+0) = 2.
\]
这意味着 \(\eta\) 是两个互不相同的不可约特征标之和，且每个重数为 1. 设 \(\chi_1\) 是 \(S_3\) 的主特征标，即对所有 \(\sigma \in S_3\) 有 \(\chi_1(\sigma) = 1\). 则
\[
(\eta, \chi_1) = \frac{1}{6}\left[1\cdot\eta(1)\chi_1(1)+3\cdot\eta((1\ 2))\chi_1((1\ 2))+2\cdot\eta((1\ 2\ 3))\chi_1((1\ 2\ 3))\right] = \frac{1}{6}(3+3+0) = 1,
\]
故主特征标作为 \(\eta\) 的成分出现，重数为 1. 这证明 \(\eta = \chi_1+\chi_2\)，其中 \(\chi_2\) 是 \(S_3\) 的某个 2 次不可约特征标（这与我们前面关于该表示的分解一致）。这也表明 \(\chi_2\) 在 \(\sigma \in S_3\) 上的值是不动点个数减 1.

（2）设 \(G = S_4\)，\(\eta\) 是 4 次自然置换特征标（于是同样 \(\eta(\sigma)\) 是不动点个数）。\(S_4\) 的共轭类由 \(1\)、\((1\ 2)\)、\((1\ 2\ 3)\)、\((1\ 2\ 3\ 4)\) 与 \((1\ 2)(3\ 4)\) 代表，大小分别为 1、6、8、6、3. 再次计算：
\[
\|\eta\|^2 = \frac{1}{24}\left[1\cdot\eta(1)^2+6\cdot\eta((1\ 2))^2+8\cdot\eta((1\ 2\ 3))^2+6\cdot\eta((1\ 2\ 3\ 4))^2+3\cdot\eta((1\ 2)(3\ 4))^2\right]
\]
\[
= \frac{1}{24}(16+24+8+0+0) = 2,
\]
故 \(\eta\) 有两个互不相同的不可约成分。若 \(\chi_1\) 是 \(S_4\) 的主特征标，则
\[
(\eta, \chi_1) = \frac{1}{24}\left[1\cdot\eta(1)+6\cdot\eta((1\ 2))+8\cdot\eta((1\ 2\ 3))+6\cdot\eta((1\ 2\ 3\ 4))+3\cdot\eta((1\ 2)(3\ 4))\right]
\]
\[
= \frac{1}{24}(4+12+8+0+0) = 1.
\]
这证明 4 次置换特征标是主特征标与一个 3 次不可约特征标之和。

（3）设 \(G = D_8\)，其中
\[
D_8 = \langle r, s \mid s^2 = r^4 = 1,\ rs = sr^{-1}\rangle.
\]
\(D_8\) 的共轭类由 \(1\)、\(s\)、\(r\)、\(r^2\) 与 \(sr\) 代表，大小分别为 1、2、2、1、2. 设 \(\varphi\) 是第 1 节例 6 中把正方形嵌入 \(\mathbb{R}^2\) 所得的 \(D_8\) 的 2 次矩阵表示：
\[
\varphi(s) = \begin{pmatrix}0 & 1\\ 1 & 0\end{pmatrix}, \quad \varphi(r) = \begin{pmatrix}0 & -1\\ 1 & 0\end{pmatrix}, \quad \varphi(r^2) = \begin{pmatrix}-1 & 0\\ 0 & -1\end{pmatrix}, \quad \varphi(sr) = \begin{pmatrix}1 & 0\\ 0 & -1\end{pmatrix}.
\]
设 \(\psi\) 是这个表示的特征标（我们把实矩阵看作复矩阵的子集）。同样，由于 \(\psi\) 实值，计算
\[
\|\psi\|^2 = \frac{1}{8}\left[1\cdot\psi(1)^2+2\psi(s)^2+2\psi(r)^2+1\cdot\psi(r^2)^2+2\psi(sr)^2\right] = \frac{1}{8}(4+0+0+4+0) = 1.
\]
这证明表示 \(\varphi\) 不可约（即使允许用复矩阵作相似变换）。

我们已经看到两个特征标之和仍是特征标。具体地，若 \(\psi_1\) 与 \(\psi_2\) 是表示 \(\varphi_1\) 与 \(\varphi_2\) 的特征标，则 \(\psi_1+\psi_2\) 是 \(\varphi_1+\varphi_2\) 的特征标。

\begin{proposition}\label{pro:18.17}
若 \(\psi_1\) 与 \(\psi_2\) 是特征标，则它们的乘积 \(\psi_1\psi_2\) 也是特征标。
\end{proposition}

\begin{proof}
设 \(V_1\) 与 \(V_2\) 是提供特征标 \(\psi_1\) 与 \(\psi_2\) 的 \(\mathbb{C}G\)-模，定义 \(W = V_1 \otimes_{\mathbb{C}}V_2\). 由于每个 \(g \in G\) 作为线性变换作用在 \(V_1\) 与 \(V_2\) 上，\(g\) 在简单张量上由 \(g(v_1 \otimes v_2) = (gv_1)\otimes(gv_2)\) 定义的作用按线性性延拓为 \(W\) 上良定义的线性变换（第 11.2 节命题 17）。容易验证这个作用也使 \(W\) 成为 \(\mathbb{C}G\)-模。由第 11.2 节习题 38，\(W\) 提供的特征标是 \(\psi_1\psi_2\).
\end{proof}

下一章将包含进一步的特征标显式计算，以及群特征标在证明某些类群的定理时的应用。

\subsection*{关于傅里叶分析与群特征标的一些注记}

这段简短讨论意在指出上述结果与数学其他领域的一些联系。

至此描述的群表示论是称为\textbf{调和分析}的数学领域的一个特殊分支。读者可能已经熟悉也属于这一领域的 Fourier 级数基本理论。我们作一些观察，说明有限群的表示论如何对应于某些无限群的「Fourier 级数」（特别地，对应于圆上的 Fourier 级数）。为了数学上精确，需要 Lebesgue 积分来保证某些（Hilbert）空间的完备性，但读者把「Lebesgue」换成「Riemann」也能体会其中的意味。

设 \(G\) 是 \(\mathbb{C}\) 中单位圆上点的乘法群：
\[
G = \{z \in \mathbb{C} \mid |z| = 1\}.
\]
我们通常把 \(G\) 看作 \(\mathbb{R}\) 中区间 \([0, 2\pi]\) 且两个端点等同，即看作加法群 \(\mathbb{R}/2\pi\mathbb{Z}\)（同构是：实数 \(x\) 对应复数 \(e^{ix}\)）。注意 \(G\) 有平移不变的测度，即 Lebesgue 测度，而圆的测度为 \(2\pi\). 对有限群，计数测度是平移不变的测度（故子集的测度是该子集中元素的个数），而有限群上关于这个计数测度的积分就是有限和。

空间
\[
L^2(G) = \{f : G \to \mathbb{C} \mid f \text{ 可测且 } |f|^2 \text{ 在 } G \text{ 上可积}\}
\]
扮演无限群 \(G\) 的群代数的角色。这个空间在函数的\textbf{卷积}下成为含 1 的交换环：对 \(f, g \in L^2(G)\)，乘积 \(f * g : G \to \mathbb{C}\) 定义为
\[
(f*g)(x) = \frac{1}{2\pi}\int_0^{2\pi}f(x-y)g(y)\,dy \qquad \text{对所有 } x \in G.
\]
（回忆对有限群 \(H\)，群代数形式上也是 \(H\) 上 \(\mathbb{C}\)-值函数之环，其乘法为卷积，且这些函数写成形式和——元素 \(\sum a_gg \in \mathbb{C}G\) 表示把 \(g\) 映到 \(a_g \in \mathbb{C}\) 的函数。）

\(G\) 到 \(GL_1(\mathbb{C})\) 的连续同态的完全集合由
\[
e_n(x) = e^{inx}, \qquad x \in [0, 2\pi],\quad n \in \mathbb{Z}
\]
给出。（回忆对有限阿贝尔群，所有不可约表示都是 1 维的，而对 1 维表示，特征标与表示可以等同。）

环 \(L^2(G)\) 具有 Hermite 内积：对 \(f, g \in L^2(G)\)，
\[
(f, g) = \frac{1}{2\pi}\int_0^{2\pi}f(t)\overline{g(t)}\,dt.
\]
关于这个内积，\(\{e_n \mid n \in \mathbb{Z}\}\) 是一组标准正交基（其中「基」一词用于分析意义：这些元素独立，且与所有它们正交的唯一函数是 0）。此外，
\[
L^2(G) = \widehat{\bigoplus}_{n \in \mathbb{Z}}E_n,
\]
其中 \(E_n\) 是由 \(e_n\) 张成的 1 维子空间，直和上的帽子表示取 \(L^2\)-拓扑下的闭包，等号表示 \(L^2\) 意义下的相等。（回忆有限阿贝尔群的群代数是不可约 1 维子模的直和，每个重数为 1。）这些事实蕴含 Fourier 分析中的著名结果：\([0, 2\pi]\) 上每个平方可积函数 \(f(x)\) 都有 Fourier 级数
\[
f(x) \sim \sum_{n=-\infty}^{\infty}c_ne^{inx},
\]
其中 Fourier 系数 \(c_n\) 由
\[
c_n = (f, e_n) = \frac{1}{2\pi}\int_0^{2\pi}f(t)e^{-int}\,dt
\]
给出。

这段简短描述指出了有限群的表示论如何推广到某些无限群，而我们已经证明的结果在后者的语境中可能已为人所熟悉。事实上，对任意（不必阿贝尔的）紧 Lie 群有完全类似的理论——这里不可约（复）表示不必是 1 维的，但它们都是有限维的，而 \(L^2(G)\) 分解为它们的直和，每个出现的重数等于它的次数。这个理论（至少在入门层面）的重点常常在于能把函数表示为（Fourier）级数，然后利用这些级数去解决其他问题（例如解微分方程）。底层的群提供了建立这种「调和分析」所依据的「对称性」，而它本身并不是研究的主要对象。

\subsection*{习题}

设 \(G\) 是有限群。除另有说明外，所有表示与特征标都在 \(\mathbb{C}\) 上。

\begin{enumerate}
\item 证明对元素取自任意交换环的 \(n \times n\) 矩阵 \(A\) 与 \(B\) 有 \(\operatorname{tr}AB = \operatorname{tr}BA\).

\item 在（a）至（c）中，设 \(\psi\) 是指定表示 \(\varphi\) 提供的特征标。

（a）设 \(\varphi\) 是第 1 节第二组例子中例 6 描述的 \(D_{10}\) 的 2 次表示（这里 \(n = 5\)），证明 \(\|\psi\|^2 = 1\)（故 \(\varphi\) 不可约）。

（b）设 \(\varphi\) 是第 1 节第二组例子中例 7 描述的 \(Q_8\) 的 2 次表示，证明 \(\|\psi\|^2 = 1\)（故 \(\varphi\) 不可约）。

（c）设 \(\varphi\) 是第 1 节第二组例子中例 8 描述的 \(Q_8\) 的 4 次表示，证明 \(\|\psi\|^2 = 4\)（故虽然 \(\varphi\) 在 \(\mathbb{R}\) 上不可约，它在 \(\mathbb{C}\) 上分解为某个 2 次不可约表示的两倍）。

\item 若 \(\chi\) 是 \(G\) 的不可约特征标，证明 \(\mathbb{C}G\)-模的 \(\chi\)-等型子空间是唯一的。

\item 证明若 \(N\) 是任意不可约 \(\mathbb{C}G\)-模且 \(M = N \oplus N\)，则 \(M\) 有无穷多把它分解为两个 \(N\) 的拷贝的直和分解。

\item 证明类函数是特征标当且仅当它是不可约特征标的非负整数线性组合。

\item 设 \(\varphi : G \to GL(V)\) 是特征标为 \(\psi\) 的表示。设 \(W = \{v \in V \mid \varphi(g)(v) = v \text{ 对所有 } g \in G\}\) 是 \(V\) 中被 \(G\) 的所有元素逐点固定的子空间。证明 \(\dim W = (\psi, \chi_1)\)，其中 \(\chi_1\) 是 \(G\) 的主特征标。

\item 假设 \(V\) 是 \(\mathbb{C}G\)-模，\(G\) 通过置换基 \(\mathcal{B} = \{e_1, \ldots, e_n\}\) 作用。把 \(\mathcal{B}\) 写成 \(G\) 在 \(\mathcal{B}\) 上的轨道 \(\mathcal{B}_1, \ldots, \mathcal{B}_t\) 的不交并。

（a）证明作为 \(\mathbb{C}G\)-模有 \(V = V_1 \oplus \cdots \oplus V_t\)，其中 \(V_i\) 是 \(\mathcal{B}_i\) 的张成。

（b）证明若 \(v_i\) 是 \(\mathcal{B}_i\) 中向量之和，则由 \(v_i\) 张成的 \(V_i\) 的 1 维子空间是 \(V_i\) 中提供平凡表示的唯一的 \(\mathbb{C}G\)-子模（换句话说，\(V_i\) 中在 \(G\) 的作用下固定的任何向量都是 \(v_i\) 的倍数）。〔用 \(G\) 在 \(\mathcal{B}_i\) 上传递这一事实。另见第 1 节习题 8。〕

（c）设 \(W = \{v \in V \mid \varphi(g)(v) = v \text{ 对所有 } g \in G\}\) 是 \(V\) 中被 \(G\) 的所有元素逐点固定的子空间。推出 \(\dim W = t\)，即 \(G\) 在 \(\mathcal{B}\) 上的轨道数。

\item 证明下述结果（有时称为 \textbf{Burnside 引理}，虽然它源于 Frobenius）：设 \(G\) 是 \(S_n\) 的子群，对每个 \(\sigma \in G\) 用 \(\operatorname{Fix}(\sigma)\) 表示 \(\sigma\) 在 \(\{1, \ldots, n\}\) 上的不动点个数。设 \(t\) 是 \(G\) 在 \(\{1, \ldots, n\}\) 上的轨道数。则
\[
t|G| = \sum_{g \in G}\operatorname{Fix}(g).
\]
〔用上面两个习题。〕

\item 设 \(G\) 是有限集 \(\Omega\) 上的非平凡传递置换群，\(\psi\) 是由 \(\Omega\) 上的置换（参看第 1 节例 4）所得的 \(\mathbb{C}\) 上线性表示提供的特征标，故 \(\psi(\sigma)\) 是 \(\sigma\) 在 \(\Omega\) 上的不动点个数。现在设 \(G\) 通过 \(g \cdot (\omega_1, \omega_2) = (g \cdot \omega_1, g \cdot \omega_2)\) 作用在集合 \(\Omega \times \Omega\) 上，并设 \(\pi\) 是这个作用所得的线性表示提供的特征标。

（a）证明 \(\pi = \psi^2\).

（b）证明 \(G\) 在 \(\Omega \times \Omega\) 上的轨道数由内积 \((\psi, \psi)\) 给出。〔由上面几个习题，\(\Omega \times \Omega\) 上的轨道数等于 \((\pi, \chi_1)\)，其中 \(\chi_1\) 是主特征标。〕

（c）回忆若 \(G\) 在 \(\Omega \times \Omega\) 上的作用恰好有 2 个轨道，则称它在 \(\Omega\) 上\textbf{二重传递}（它总有至少 2 个轨道，因为对角线 \(\{(\omega, \omega) \mid \omega \in \Omega\}\) 是一个轨道）。证明若 \(G\) 在 \(\Omega\) 上二重传递，则 \(\psi = \chi_1+\chi_2\)，其中 \(\chi_1\) 是主特征标而 \(\chi_2\) 是 \(G\) 的非主不可约特征标。

（d）设 \(\Omega = \{1, 2, \ldots, n\}\)，\(G = S_n\) 按自然方式作用在 \(\Omega\) 上。证明相应线性表示的特征标分解为主特征标加上一个 \(n-1\) 次不可约特征标。

\item 设 \(\psi\) 是群 \(G\) 的任意 2 次表示的特征标，\(x\) 是 \(G\) 中一个 2 阶元素。证明 \(\psi(x) = 2\)、0 或 \(-2\). 把这个结果推广到 \(n\) 次表示。

\item 设 \(\chi\) 是 \(G\) 的不可约特征标。证明对 \(G\) 的中心中每个元素 \(z\) 有 \(\chi(z) = \varepsilon\chi(1)\)，其中 \(\varepsilon\) 是 \(\mathbb{C}\) 中某个单位根。〔用 Schur 引理。〕

\item 设 \(\psi\) 是 \(G\) 的某个表示 \(\varphi\) 的特征标。证明对 \(g \in G\) 有：

（a）若 \(\psi(g) = \psi(1)\)，则 \(g \in \ker\varphi\)；以及

（b）若 \(|\psi(g)| = \psi(1)\) 且 \(\varphi\) 忠实，则 \(g \in Z(G)\)（其中 \(|\psi(g)|\) 是 \(\psi(g)\) 的复绝对值）。〔用命题 14 的证明方法。〕

\item 设 \(\varphi : G \to GL(V)\) 是表示，\(\chi : G \to \mathbb{C}^\times\) 是 1 次表示。证明由 \(\chi\varphi(g) = \chi(g)\varphi(g)\) 定义的 \(\chi\varphi : G \to GL(V)\) 是表示（注意线性变换 \(\varphi(g)\) 与复数 \(\chi(g)\) 的乘法是良定义的）。证明 \(\chi\varphi\) 不可约当且仅当 \(\varphi\) 不可约。证明若 \(\psi\) 是 \(\varphi\) 提供的特征标，则 \(\chi\psi\) 是 \(\chi\varphi\) 提供的特征标。推出：任意不可约特征标与任意 1 次特征标的乘积仍不可约。
\end{enumerate}

下面几个习题研究代数共轭特征标的概念。这些习题可以看作命题 14 的推广及这些推广的一些推论。特别地，我们得到「群的所有不可约特征标都取有理值」这一条件的群论刻画。

设 \(F\) 是 \(\mathbb{C}\) 中所有在 \(\mathbb{Q}\) 上代数的元素构成的子域（代数数域）。故 \(F\) 是 \(\mathbb{C}\) 中包含的 \(\mathbb{Q}\) 的代数闭包，且前面在 \(\mathbb{C}\) 上建立的所有结果不加改变地在 \(F\) 上成立。

\begin{enumerate}
\setcounter{enumi}{13}
\item 注意由于 \(F \subseteq \mathbb{C}\)，每个表示 \(\varphi : G \to GL_m(F)\) 也可看作复表示。证明若 \(\varphi\) 是 \(F\) 上的表示且在 \(F\) 上不可约，则把它看作更大域 \(\mathbb{C}\) 上的表示时也不可约（注意若 \(F\) 不代数闭，这并不成立——参看上面习题 2（c））。证明 \(G\) 在 \(F\) 上的不可约特征标集与在 \(\mathbb{C}\) 上的不可约特征标集相同（即它们是 \(G\) 上完全相同的一组类函数）。推出每个复表示都等价于 \(F\) 上的一个表示。〔由于 \(F\) 是特征为 0 的代数闭域，\(F\) 或 \(\mathbb{C}\) 上的不可约特征标都由第一正交关系刻画。〕
\end{enumerate}

设 \(\varphi : G \to GL_m(F)\) 是任意表示，特征标为 \(\psi\). 用 \(\mathbb{Q}(\varphi)\) 表示由所有 \(g \in G\) 的矩阵 \(\varphi(g)\) 的全部元素生成的 \(F\) 的子域。

\begin{enumerate}
\setcounter{enumi}{14}
\item 证明 \(\mathbb{Q}(\varphi)\) 是 \(\mathbb{Q}\) 的有限扩张。
\end{enumerate}

现在设 \(K\) 是任何包含 \(\mathbb{Q}(\varphi)\) 的 \(\mathbb{Q}\) 的伽罗瓦扩张，\(\sigma \in \operatorname{Gal}(K/\mathbb{Q})\). 事实上，由于 \(K\) 的每个自同构都可延拓为 \(F\) 的自同构，我们可以假定 \(\sigma\) 是 \(F\) 的任意自同构。映射 \(\varphi^\sigma : G \to GL_n(F)\) 定义如下：\(\varphi^\sigma(g)\) 是把域自同构 \(\sigma\) 作用于矩阵 \(\varphi(g)\) 的元素所得的元素构成的 \(n \times n\) 矩阵。

\begin{enumerate}
\setcounter{enumi}{15}
\item 证明 \(\varphi^\sigma\) 是表示。并证明 \(\varphi^\sigma\) 的特征标是 \(\psi^\sigma\)，其中 \(\psi^\sigma(g) = \sigma(\psi(g))\).

\item 证明 \(\varphi\) 不可约当且仅当 \(\varphi^\sigma\) 不可约。
\end{enumerate}

表示 \(\varphi^\sigma\)（或特征标 \(\psi^\sigma\)）称为 \(\varphi\)（或 \(\psi\)）在 \(\sigma\) 下的\textbf{代数共轭}；若存在 \(F\) 的某个自同构 \(\sigma\) 使 \(\varphi_1^\sigma = \varphi_2\)（或 \(\psi_1^\sigma = \psi_2\)），则称两个表示 \(\varphi_1\) 与 \(\varphi_2\)（或特征标 \(\psi_1\) 与 \(\psi_2\)）\textbf{代数共轭}。对这种（标准）记号需要小心，因为指数记号通常表示右作用，而 \(F\) 的自同构在表示上左作用：\(\varphi^{\sigma\tau} = (\varphi^\tau)^\sigma\).

设 \(\mathbb{Q}(\psi)\) 是由所有 \(g \in G\) 的数 \(\psi(g)\) 生成的 \(F\) 的子域。设 \(|G| = n\)，\(\varepsilon\) 是 \(F\) 中的本原 \(n\) 次单位根。

\begin{enumerate}
\setcounter{enumi}{17}
\item 证明 \(\mathbb{Q}(\psi) \subseteq \mathbb{Q}(\varepsilon)\). 推出 \(\mathbb{Q}(\psi)\) 是 \(\mathbb{Q}\) 的伽罗瓦扩张且伽罗瓦群阿贝尔。〔见命题 14。〕
\end{enumerate}

回忆第 14.5 节中 \(\operatorname{Gal}(\mathbb{Q}(\varepsilon)/\mathbb{Q}) \cong (\mathbb{Z}/n\mathbb{Z})^\times\)，其中伽罗瓦自同构由它们在本原元 \(\varepsilon\) 上的作用给出：\(\sigma_a : \varepsilon \mapsto \varepsilon^a\)，其中 \(a\) 是与 \(n\) 互素的整数。

\begin{enumerate}
\setcounter{enumi}{18}
\item 证明若 \(\sigma_a \in \operatorname{Gal}(\mathbb{Q}(\varepsilon)/\mathbb{Q})\) 是上述定义的域自同构，则对所有 \(g \in G\) 有 \(\psi^{\sigma_a}(g) = \psi(g^a)\).〔用命题 14 的方法。〕

\item 证明若 \(g\) 是 \(G\) 中与 \(g^a\) 共轭的元素（对所有与 \(n = |G|\) 互素的整数 \(a\)），则对 \(G\) 的每个特征标 \(\psi\) 有 \(\psi(g) \in \mathbb{Q}\).〔用上一习题以及「\(\mathbb{Q}\) 是被所有 \(\sigma_a\) 固定的域」这一事实。〕

\item 证明元素 \(g \in G\) 对所有与 \(|G|\) 互素的整数 \(a\) 都与 \(g^a\) 共轭，当且仅当它对所有与 \(|g|\) 互素的整数 \(a'\) 都与 \(g^{a'}\) 共轭。

\item 证明对任何正整数 \(n\)，对称群 \(S_n\) 的每个特征标都是有理值的（即对所有 \(g \in S_n\) 与所有特征标 \(\psi\) 有 \(\psi(g) \in \mathbb{Q}\)）。
\end{enumerate}

下面两个习题建立习题 20 的逆。

\begin{enumerate}
\setcounter{enumi}{22}
\item 证明群 \(G\) 中元素 \(x\) 与 \(y\) 共轭当且仅当对 \(G\) 的所有不可约特征标 \(\chi\) 有 \(\chi(x) = \chi(y)\).

\item 设 \(g \in G\)，并假设 \(G\) 的每个不可约特征标在 \(g\) 上都有理值。证明对每个与 \(|G|\) 互素的整数 \(a\)，\(g\) 与 \(g^a\) 共轭。〔若对某个与 \(|G|\) 互素的 \(a\)，\(g\) 不与 \(g^a\) 共轭，则由上一习题存在不可约特征标 \(\chi\) 使 \(\chi(g) \neq \chi(g^a)\). 由假设 \(\chi(g) \in \mathbb{Q}\) 导出矛盾。〕

\item 描述 \(n\) 阶循环群的哪些不可约特征标彼此代数共轭。

\item 证明 \(Q_8\) 与 \(D_8\) 的每个不可约特征标都有理值。证明 \(D_{10}\) 有一个不可约特征标不是有理值的。

\item 设 \(G = H \times K\)，\(\varphi : H \to GL(V)\) 是 \(H\) 的不可约表示，特征标为 \(\chi\). 则 \(G \to H \to GL(V)\) 给出 \(G\) 的不可约表示，其中第一个映射是自然投射；这个表示的特征标 \(\tilde\chi\) 满足 \(\tilde\chi((h, k)) = \chi(h)\). 同样，\(K\) 的任何不可约特征标 \(\psi\) 给出 \(G\) 的不可约特征标 \(\tilde\psi\)，满足 \(\tilde\psi((h, k)) = \psi(k)\).

（a）证明乘积 \(\tilde\chi\tilde\psi\) 是 \(G\) 的不可约特征标。〔证明它的范数为 1。〕

（b）证明 \(G\) 的每个不可约特征标都由诸直因子的不可约特征标的这种乘积得到。〔用定理 10 的（3）或（4）。〕

\item （\(\operatorname{GL}_2(\mathbb{Q})\) 的有限子群）设 \(G\) 是 \(\operatorname{GL}_2(\mathbb{Q})\) 的有限子群。

（a）证明 \(\operatorname{GL}_2(\mathbb{Q})\) 不含 \(n = 5\)、7 或 \(n \geq 9\) 阶元素。推出 \(|G| = 2^a3^b\).〔用有理标准形。〕

（b）证明 Klein 四元群是 \(\operatorname{GL}_2(\mathbb{Q})\) 中唯一的非循环阿贝尔子群。由（a）推出 \(|G| \leq 24\).

（c）证明 \(\operatorname{GL}_2(\mathbb{Q})\) 的有限子群只有 1、2、3、4、6 阶循环群、Klein 四元群以及 6、8、12 阶二面体群。〔用第 4.5 节小阶群的分类与第 1 节习题 10 把 \(G\) 限制到这个清单。反过来证明所列每个群都有 2 维忠实有理表示。〕
\end{enumerate}
